\ifx\danslelivre\undefined
\RequirePackage{luatex85}
\documentclass[10pt]{../fichiers-configuration-latex/smfart}
\usepackage{polyglossia}
%\usepackage[francais]{babel}
\usepackage{fontspec}
\usepackage[math-style=ISO,bold-style=upright,sans-style=upright]{unicode-math}
%\usepackage[mathletters]{ucs}
%\usepackage[utf8]{inputenc}
%\usepackage[T1]{fontenc}
\usepackage[unicode,colorlinks=true,linkcolor=blue,citecolor=Sepia,urlcolor=magenta]{hyperref}
\setmainlanguage{french}
\setmainfont[Ligatures=TeX]{Linux Libertine O}
% variante possible : TeXGyrePagella ?
\setmathfont{XITS Math}

% maths
\usepackage{amsmath}
%\usepackage{amssymb}
%\usepackage{mathrsfs}
%\usepackage{stmaryrd}
%\usepackage{euscript}

% couleurs
\usepackage[usenames,dvipsnames]{xcolor}

% graphiques
\usepackage{graphicx}
\usepackage[all]{xy}
% Xy-pic 3.8.9's PDF driver assumes that \fontdimen8\textfont3 is a
% legacy TeX rule width.  Under LuaTeX with unicode-math/XITS Math it is
% instead 6.62pt, which turns every arrow shaft into a thick black bar.
% Override only xypdf's private cached width; do not alter the math font.
\makeatletter
\AtBeginDocument{%
  \@ifundefined{xP@preclw}{}{%
    \edef\xP@preclw{0.4pt}%
    \edef\xP@lw{\xP@dim\xP@preclw}%
    \PackageInfo{sga4-xypdf-compat}{Xy-pic PDF line width set to \xP@preclw}%
  }%
}
\makeatother
%\usepackage{tikz}
%\usetikzlibrary{matrix,arrows,calc}
%\tikzset{column sep=3em,row sep=2em,text height=1.5ex, text depth=0.25ex,>=angle 90}
%\tikzset{description/.style={fill=white,inner sep=1pt}}

% hyperliens, 'etiquettes
\usepackage{xr-hyper}
\usepackage{../fichiers-configuration-latex/smfthm}
%\ifx\cacheretiquettes\undefined
%\usepackage[inline]{showlabels} \else \fi % pour rendre visibles les 'etiquettes
%\usepackage{../fichiers-configuration-latex/makeidx} % pour faire un index
%\usepackage[pagebackref,unicode]{hyperref}  % pour voir o`u sont cit'ees les r'ef'erences
\hypersetup{unicode,colorlinks=true,linkcolor=blue,citecolor=Sepia,urlcolor=magenta,
  hypertexnames=false}
%\hypersetup{breaklinks=true} %permet le retour `a la ligne dans les liens trop longs. Incompatible avec pdflatex
%\renewcommand*{\backref}[1]{{\small [cit\'e page(s)~#1.]}}

\usepackage{xspace}
\usepackage{pifont} % pour expose XVII

\renewcommand{\thefootnote}{\roman{footnote}} % pour mettre les notes en bas de page sous la forme (ii) et non (2), qui ressemble trop (?) a un exposant.

\newcounter{pageoriginal}
\newcommand{\pageoriginale}{\refstepcounter{pageoriginal}\marginpar{\footnotesize\textbf{\thepageoriginal}}\xspace}
\let\pageoriginaled\pageoriginale

\usepackage{ifthen}
\newboolean{original}
\def\sisi{\ifthenelse{\boolean{original}}}
\setboolean{original}{false}

\newcounter{cnote}
\def\nde#1{\ifthenelse{\boolean{original}}{}{\stepcounter{cnote}\def\thefootnote{\arabic{cnote}}\protect\footnote[\value{cnote}]{N.D.E.: #1}}}
\def\ndemark{\ifthenelse{\boolean{original}}{}{\stepcounter{cnote}\def\thefootnote{\arabic{cnote}}\protect\footnotemark[\value{cnote}]}}
\def\ndetext#1{\ifthenelse{\boolean{original}}{}{\def\thefootnote{\arabic{cnote}}\protect\footnotetext[\value{cnote}]{N.D.E.: #1}}}

\setmainlanguage{french}
\setotherlanguage{english}

% unicode-math deliberately disables the legacy \Arrowvert glyph.  The
% source uses it once as a rotated equality marker; \Vert has the intended
% appearance and keeps the historical diagram buildable with LuaLaTeX.
\renewcommand{\Arrowvert}{\Vert}


%%
%% Pour afficher le numéro de version Git (exige que LuaTeX soit appelé avec --shell-escape)
%%

\catcode`\%=12
\directlua{
function command_output(cmd)
  local output_pipe = io.popen(cmd)
  local line_iter = output_pipe and output_pipe:lines()
  local line = line_iter and line_iter()
  if output_pipe then output_pipe:close() end
  return line
end
local git_command1 = [[git log --pretty=format:"%h %ai" -1]]
local git_command2 = [[if git diff --quiet HEAD ; then echo clean ; else echo dirty ; fi]]
function git_status()
  local status = command_output(git_command1)
  if status then
    tex.tprint({luatexbase.catcodetables.CatcodeTableLaTeX, string.char(0x5c).."texttt{"}, {-2,status}, {luatexbase.catcodetables.CatcodeTableLaTeX, "}"})
    if not (command_output(git_command2) == "clean") then tex.sprint(-2, " (modifications non validées)") end
  else
    tex.sprint(-2, "?")
  end
end
}
\catcode`\%=14
\def\showgitstatus{\directlua{git_status()}}


% macros pour LaTeX
\renewcommand{\mathsf}{}

% Police différente pour U+2113 SCRIPT SMALL L
\setmathfont[range="2113]{latinmodern-math.otf}

% Give every explicit SGA label its own destination and preserve the printed
% number supplied as the optional argument.  The former definition inherited
% whichever counter anchor happened to precede it and ignored #1 entirely.
\makeatletter
\def\etiquette[#1]#2{%
  \phantomsection
  \protected@edef\@currentlabel{#1}%
  \label{#2}%
}
\makeatother
\def\eqhyperref[#1]#2{(\hyperref[#1]{#2})}
\def\pagehyperref[#1]#2{(page \hyperref[#1]{#2})}
\def\hyperrefext[#1][#2]#3{\hyperref[#2]{#1 #3}}
\def\eqhyperrefext[#1][#2]#3{\hyperref[#2]{#1 #3}}     %cf. exposes 17
\def\neweqhyperrefext[#1][#2]#3{\hyperref[#2]{#1 #3}}  %cf. exposes 17

\def\shat#1{\protect\widehat{#1}}
\def\stilde#1{\protect\widetilde{#1}}
%%%%%%%%%%%%%%%%%%


\def\etcn{etc.\xspace}
\def\etco{etc...\xspace}
\def\etc{\sisi{\etco}{\etcn}}
\def\ldotsn{\sisi{...}{\etcn}}

\def\cf{{cf.}\kern.3em}
\def\Cf{{Cf.}\kern.3em}
\def\ie{{i.e.}\kern.3em}
\def\iev{{i.e.},\xspace}
\def\loccito{loc.\kern2pt cit.\xspace}
\def\loccitn{\textit{loc.\kern2pt cit.}\xspace}
\def\loccit{\sisi{\loccito}{\loccitn}}
\def\EGA{\textup{EGA}\kern.2em}
\def\SGA{\textup{SGA}\kern.3em}

\def\quoi{}

\def\quav{\quad,\quad}
\def\quaet{\quad \text{et}\quad}
\def\fibreproduct#1#2#3{#1 ×_{#3} #2}
\let\fprod\fibreproduct
\def\ccheck#1#2{\mbox{${\check{\check{\rm #1}}}^{\lower2pt\hbox{$\scriptstyle #2$}}$}}
\def\sheaf#1{\protect\mathcal{#1}}
\def\vieuxsheaf#1{{#1}} % cf. problème avec sheaf{a\Phi_i} par exemple...
\DeclareMathOperator{\Hom}{Hom}
\DeclareMathOperator{\can}{can}
\DeclareMathOperator{\Ind}{Ind}
\DeclareMathOperator{\ind}{ind}
%\DeclareMathOperator{\Hhomtop}{\mathcal{H}\!\mathit{omtop}}
\DeclareMathOperator{\Hhomtop}{\mathbf{Homtop}}
\DeclareMathOperator{\Point}{Point}
\DeclareMathOperator{\Pro}{Pro}
\DeclareMathOperator{\Top}{Top}
\DeclareMathOperator{\Mod}{Mod}
\DeclareMathOperator{\Cart}{Cart}
%\DeclareMathOperator{\Mmorsite}{\mathcal{M}\!\mathit{orsite}}
\DeclareMathOperator{\Mmorsite}{\mathbf{Morsite}}
\DeclareMathOperator{\loc}{loc}
\DeclareMathOperator{\Loc}{Loc}
\DeclareMathOperator{\Pic}{Pic}
\DeclareMathOperator{\Et}{Et}
\DeclareMathOperator{\card}{card}
\DeclareMathOperator{\Ens}{Ens}
\DeclareMathOperator{\h}{\underline{H}}
\DeclareMathOperator{\Fonct}{Fonct}
\DeclareMathOperator{\Ker}{Ker}
\DeclareMathOperator{\but}{but}
\DeclareMathOperator{\Ob}{Ob}
\DeclareMathOperator{\ob}{ob}
\DeclareMathOperator{\Lib}{Lib}
\DeclareMathOperator{\esj}{esj}
\DeclareMathOperator{\id}{id}
\DeclareMathOperator{\Id}{Id}
\DeclareMathOperator{\ad}{ad}
\DeclareMathOperator{\Iim}{Im}
\DeclareMathOperator{\pr}{pr}
\DeclareMathOperator{\Cat}{Cat}
\DeclareMathOperator{\car}{car}
\DeclareMathOperator{\SysInd}{SysInd}
\DeclareMathOperator{\Sup}{Sup}
\DeclareMathOperator{\Coker}{Coker}
\DeclareMathOperator{\Filt}{Filt}
\DeclareMathOperator{\acc}{acc}
\DeclareMathOperator{\sur}{sur}
\DeclareMathOperator{\cat}{cat}
\DeclareMathOperator{\Ord}{Ord}
\DeclareMathOperator{\Gr}{Gr}
\DeclareMathOperator{\Cov}{Cov}
\DeclareMathOperator{\Caf}{Caf}
\DeclareMathOperator{\coker}{coker}
\DeclareMathOperator{\coim}{coim}
\DeclareMathOperator{\im}{im}
\DeclareMathOperator{\Ouv}{Ouv}
\DeclareMathOperator{\diag}{diag}
\DeclareMathOperator{\Esp}{Esp}
\DeclareMathOperator{\TOP}{TOP}
\DeclareMathOperator{\Unif}{Unif}
\DeclareMathOperator{\Iind}{\mathbf{Ind}}
\DeclareMathOperator{\Hhomcart}{\mathbf{Homcart}}
\DeclareMathOperator{\Hhomtopess}{\mathbf{Homtopess}}
\DeclareMathOperator{\Ppoints}{\mathbf{Points}}
\DeclareMathOperator{\Ppoint}{\mathbf{Point}}
\DeclareMathOperator{\Ttors}{\mathbf{Tors}}
\DeclareMathOperator{\Ssex}{\mathbf{Sex}}
\DeclareMathOperator{\Ffib}{\mathbf{Fib}}
\DeclareMathOperator{\Vvois}{\mathbf{Vois}}
\DeclareMathOperator{\Ppointess}{\mathbf{Pointess}}
\DeclareMathOperator{\scrHom}{\mathbf{Hom}}
\DeclareMathOperator{\sheafPic}{\mathbf{Pic}}
\DeclareMathOperator{\sheafAut}{\mathbf{Aut}}
\DeclareMathOperator{\Hhomtoplocan}{\mathbf{Homtoplocan}}
\DeclareMathOperator{\Hhomtopan}{\mathbf{Homtopan}}
\DeclareMathOperator{\Pptgeom}{\mathbf{Ptgéom}}
\DeclareMathOperator{\Morsite}{Morsite}
\DeclareMathOperator{\Fib}{Fib}
\DeclareMathOperator{\Fais}{Fais}
\DeclareMathOperator{\kar}{kar}
\DeclareMathOperator{\Groupes}{Groupes}
\DeclareMathOperator{\Restr}{Restr}
\DeclareMathOperator{\Isom}{Isom}
\DeclareMathOperator{\Prol}{Prol}
\DeclareMathOperator{\Res}{Res}
\DeclareMathOperator{\Tors}{Tors}
\DeclareMathOperator{\esp}{esp}
\DeclareMathOperator{\es}{es}
\DeclareMathOperator{\od}{od}
\DeclareMathOperator{\gr}{gr}
\DeclareMathOperator{\sob}{sob}
\DeclareMathOperator{\dfn}{dfn}
\DeclareMathOperator{\Sch}{Sch}
\DeclareMathOperator{\ouv}{ouv}
\DeclareMathOperator{\cons}{cons}
\DeclareMathOperator{\Nerf}{Nerf}
\DeclareMathOperator{\Points}{Points}
\DeclareMathOperator{\Espsob}{Espsob}
\DeclareMathOperator{\Homtop}{Homtop}
\DeclareMathOperator{\Pt}{Pt}
\DeclareMathOperator{\Quot}{Quot}
\DeclareMathOperator{\glob}{glob}
\DeclareMathOperator{\Homfib}{Homfib}
\DeclareMathOperator{\Qouv}{Qouv}
\DeclareMathOperator{\Quotouv}{Quotouv}
\DeclareMathOperator{\ens}{ens}
\DeclareMathOperator{\ab}{ab}
\DeclareMathOperator{\an}{an}
\DeclareMathOperator{\fs}{fs}
\DeclareMathOperator{\red}{red}
\DeclareMathOperator{\Ab}{Ab}
\DeclareMathOperator{\cosq}{cosq}
\DeclareMathOperator{\Ann}{Ann}
\DeclareMathOperator{\Zar}{Zar}


%%%%%%%%%%%%
%%%%% SGA4-2
\DeclareMathOperator{\Sex}{Sex}
\DeclareMathOperator{\vois}{vois}
\DeclareMathOperator{\Vois}{Vois}
\DeclareMathOperator{\Ext}{Ext}
\DeclareMathOperator{\ext}{ext}
\DeclareMathOperator{\Fer}{Fer}
\DeclareMathOperator{\Add}{Add}
\DeclareMathOperator{\Cocart}{Cocart}
\DeclareMathOperator{\cosk}{cosk}
\DeclareMathOperator{\plat}{plat}

\DeclareMathOperator{\faisc}{faisc}
\DeclareMathOperator{\sk}{sk}
\DeclareMathOperator{\ch}{ch}

\def\HH{\mathbf{H}}
\def\qc{\mathrm{qc}}
\def\et{\mathrm{et}}
\def\coh{\mathrm{coh}}

\def\Hvr#1{%
\begin{array}[b]{l}
{\kern2pt\scriptstyle r}\\[-8pt]
{\kern1pt\scriptstyle \vee}\\[-5pt]
H^{#1}
\end{array}}

\def\Hvv#1{%
\begin{array}[b]{l}
\kern2pt\rotatebox{90}{${\scriptscriptstyle\ll}$}\\[-6pt]
H^{#1}
\end{array}}

\def\Hvi#1{%
\begin{array}[b]{l}
{\kern0.2pt\scriptscriptstyle \infty}\\[-8pt]
{\kern1pt\scriptstyle \vee}\\[-5pt]
H^{#1}
\end{array}}

\def\Hv#1{%
\begin{array}[b]{l}
{\kern1pt\scriptstyle \vee}\\[-5pt]
H^{#1}
\end{array}}

\DeclareMathOperator{\cart}{cart}
\DeclareMathOperator{\Ffaisc}{\mathbf{Faisc}}


\def\Ev#1{%
\begin{array}[b]{l}
{\kern1pt\scriptstyle \vee}\\[-5pt]
E'_{#1}
\end{array}}

\DeclareMathOperator{\Cont}{Cont}
\DeclareMathOperator{\Ccont}{\mathscr{C}\!\mathit{ont}}
\DeclareMathOperator{\Site}{Site}
\newcommand\MM{\mathbf{M}}
\def\limtop#1{{\displaystyle{\varprojlim_{#1}\!\mbox{\rm top }}}}
\def\Limtop{\mathop{\varprojlim\!\mbox{\rm top}}}
\DeclareMathOperator{\spec}{spec}
\def\ss{\scriptstyle{\mathscr{S}}}
\DeclareMathOperator{\Gal}{Gal}
\def\smallm{\scriptstyle{\mathscr{M}}}
\DeclareMathOperator{\End}{End}
\DeclareMathOperator{\ass}{ass}
\DeclareMathOperator{\zar}{zar}
\DeclareMathOperator{\Ogg}{Ogg}
\DeclareMathOperator{\res}{res}
\DeclareMathOperator{\tr}{tr}
\DeclareMathOperator{\Supp}{Supp}
\DeclareMathOperator{\codim}{codim}
\DeclareMathOperator{\rk}{rk}
\DeclareMathOperator{\cur}{cur}
\DeclareMathOperator{\cdqc}{cdqc}
\newcommand\DD{\mathbf{D}}
\DeclareMathOperator{\cl}{cl}
\def\JJ{\mathbf{J}}
\def\FF{\mathbf{F}}
\def\BB{\mathbf{B}}
\DeclareMathOperator{\Alget}{Alget}
\DeclareMathOperator{\Ass}{Ass}
\DeclareMathOperator{\supp}{supp}
\DeclareMathOperator{\rad}{rad}
\DeclareMathOperator{\prof}{prof}
\DeclareMathOperator{\cint}{int}
\DeclareMathOperator{\sic}{sic}
\DeclareMathOperator{\Tor}{Tor}
\DeclareMathOperator{\tors}{tors}
\DeclareMathOperator{\Sym}{Sym}
\DeclareMathOperator{\parf}{parf}
\DeclareMathOperator{\tor}{tor}
\DeclareMathOperator{\defn}{defn}
\DeclareMathOperator{\Div}{Div}
\DeclareMathOperator{\Pol}{Pol}
\DeclareMathOperator{\Fex}{Fex}
\DeclareMathOperator{\Lie}{Lie}
\DeclareMathOperator{\tot}{tot}
\DeclareMathOperator{\Aut}{Aut}
\DeclareMathOperator{\HOM}{HOM}
\DeclareMathOperator{\pch}{pch}
\DeclareMathOperator{\Ch}{Ch}
\DeclareMathOperator{\EXT}{EXT}
\DeclareMathOperator{\PIC}{PIC}
\DeclareMathOperator{\Rf}{Rf}
\DeclareMathOperator{\inj}{inj}
\newcommand\VV{\mathbf{V}}

\def\Xv{%
\begin{array}[b]{l}
{\kern1pt\scriptstyle \vee}\\[-5pt]
X
\end{array}}

%%%%%%%
\newcommand\TT{\mathbf{T}}
\newcommand\XX{\mathbf{X}}
\newcommand\EE{\mathbf{E}}
%%%%%%

\def\qedsymbol{\textup{C.Q.F.D.}}
\let\CQFD\qed
\let\cqfd\qed

%%%%%\bigdot
\font\bigsymb=cmsy10 at 4pt
\def\bigdot{{\kern1.2pt\raise 1.5pt\hbox{\bigsymb\char15}}}
%%%%%

\def\Lright{\mathop{\vcenter{\hbox{\hbox{$\longrightarrow$}\lower8pt \hbox{\kern-16pt\hbox{$\longrightarrow$}}}}}}

\def\Aa{{\mathcal A}}
\def\Bb{{\mathcal B}}
\def\Cc{{\mathcal C}}
\def\Dd{{\mathcal D}}
\def\Ee{{\mathcal E}}
\def\Ff{{\mathcal F}}
\def\Gg{{\mathcal G}}
\def\Hh{{\mathcal H}}
\def\Ii{{\mathcal I}}
\def\Jj{{\mathcal J}}
\def\Kk{{\mathcal K}}
\def\Ll{{\mathcal L}}
\def\Mm{{\mathcal M}}
\def\Nn{{\mathcal N}}
\def\Oo{{\mathcal O}}
\def\Pp{{\mathcal P}}
\def\Qq{{\mathcal Q}}
\def\Rr{{\mathcal R}}
\def\Ss{{\mathcal S}}
\def\Tt{{\mathcal T}}
\def\Uu{{\mathcal U}}
\def\Vv{{\mathcal V}}
\def\Ww{{\mathcal W}}
\def\Xx{{\mathcal X}}
\def\Yy{{\mathcal Y}}
\def\Zz{{\mathcal Z}}

\newcommand{\CC}{\mathbf{C}}
\newcommand{\GG}{\mathbf{G}}
\newcommand{\KK}{\mathbf{K}}
\newcommand{\LL}{\mathbf{L}}
\newcommand{\NN}{\mathbf{N}}
\newcommand{\PP}{\mathbf{P}}
\newcommand\QQ{\mathbf{Q}}
\newcommand\RR{\mathbf{R}}
\newcommand\ZZ{\mathbf{Z}}

\def\to{\mathchoice{\longrightarrow}{\rightarrow}{\rightarrow}{\rightarrow}}
\def\from{\mathchoice{\longleftarrow}{\leftarrow}{\leftarrow}{\leftarrow}}
\def\mto{\mathchoice{\longmapsto}{\mapsto}{\mapsto}{\mapsto}}
\def\mfrom{\mathrel{\from\joinrel\mapstochar\,}}
\def\hto{\mathrel{\lhook\joinrel\to}}
\def\hfrom{\mathrel{\from\joinrel\rhook}}

\def\implies{\mathchoice{\Longrightarrow}{\Rightarrow}{\Rightarrow}{\Rightarrow}}
\def\implied{\mathchoice{\Longleftarrow}{\Leftarrow}{\Leftarrow}{\Leftarrow}}
\def\iff{\mathchoice{\Longleftrightarrow}{\Leftrightarrow}{\Leftrightarrow}{\Leftrightarrow}}

\def\To#1{\mathchoice{\xrightarrow{\labelstyle~~#1~}}{\stackrel{#1}{\longrightarrow}}{}{}}
\def\Mto#1{\mathrel{\mapstochar\To{#1}}}
\def\From#1{\mathchoice{\xleftarrow{\labelstyle~~#1~}}{\stackrel{#1}{\longleftarrow}}{}{}}

\newcommand{\isomto}{\overset{\sim}{\longrightarrow}}
\newcommand{\isomfrom}{\overset{\sim}{\longleftarrow}}

\newcommand{\IFF}{\unskip~$\iff$~\ignorespaces}
\newcommand{\THEN}{\unskip~$\implies$~\ignorespaces}

\let\leq\leqslant
\let\geq\geqslant
\let\le\leqslant
\let\ge\geqslant

\def\cf{{cf.}\kern.3em}
\def\Cf{{Cf.}\kern.3em}
\def\ie{{i.e.}\kern.3em}
\def\iev{{i.e.},\xspace}
%\def\loccito{loc.\kern2pt cit.\xspace}
%\def\loccitn{\textit{loc.\kern2pt cit.}\xspace}
%\def\loccit{\sisi{\loccito}{\loccitn}}
\def\EGA{\textup{EGA}\kern.2em}
\def\SGA{\textup{SGA}\kern.3em}

\let\oldS\S
\def\S{\oldS\kern .15em}
\def\ptbl{.\kern .2em}
\def\resp{resp.\kern.3em}
\def\Exp{Exp.\kern.3em}

\let\bar\overline
\let\hat\widehat
\let\tilde\widetilde
\let\dpl\displaystyle
\let\ldots\dots
\let\emptyset\varnothing
\let\skpt\skippointrait
\let\epsilon\varepsilon

\def\sfrac#1#2{{#1}/{#2}}

\DeclareMathOperator{\cd}{cd}
%\DeclareMathOperator{\Hhom}{\mathcal{H}\!\mathit{om}}
\DeclareMathOperator{\Hhom}{\mathbf{Hom}}
%\DeclareMathOperator{\Eext}{\mathcal{E}\!\mathit{xt}}
\DeclareMathOperator{\Eext}{\mathbf{Ext}}
\DeclareMathOperator{\Spec}{Spec}
\DeclareMathOperator{\Tr}{Tr}

\let\emptyset=\varnothing

%%%%%%%%%%%%%%%
%%% expose XVII
%%%%%%%%%%%%%%%

\makeatletter
\let\old@optexpose\optexpose
\newcommand{\empty@optexpose}[1]{}
\newcommand{\paren@optexpose}[1]{\old@optexpose{#1}(}
\newcommand{\refnoexpose}[1]{\let\optexpose\empty@optexpose\ref{#1}\let\optexpose\old@optexpose}%\refnoexpose{X.cor5.2} engendre "5.2" au lieu de "X 5.2"
\newcommand{\neweqref}[1]{\let\optexpose\paren@optexpose\textup{\ref{#1})}\let\optexpose\old@optexpose}%\neweqref{XII.eq4.1} engendre "XII (4.1)" au lieu de "(XII 4.1)"
\makeatother
\newcommand{\refvb}[2][VI${}_{\text{B}}$ ]{\sisi{#1\refnoexpose{#2}}{\ref{#2}}}

\newcommand{\leftexp}[2]{{\vphantom{#2}}^{#1}{#2}}
\newcommand{\AAA}{\mathbf{A}}
\newcommand{\Sf}{\mathfrak{S}}
\DeclareMathOperator*{\quotprojlim}{\text{``$\varprojlim$''}}
\DeclareMathOperator*{\quotinjlim}{\text{``$\varinjlim$''}}
\DeclareMathOperator{\Fl}{F\sisi{\ell}{l}}
\DeclareMathOperator{\Pl}{P\sisi{\ell}{l}}
%\DeclareMathOperator{\cl}{c\sisi{\ell}{l}}
\DeclareMathOperator{\tdf}{tdf}
\DeclareMathOperator{\rg}{rg}
\newcommand{\iInd}{\sisi{I}{i}nd}
\newcommand{\iind}{ind}
\newcommand{\ppro}{pro}
\newcommand{\tTor}{\sisi{T}{t}or}
\newcommand{\ttor}{tor}
\newcommand{\chap}{Chap.}
\newcommand{\Ffex}{\Ff\mathit{ex}}
\newcommand{\sheafk}{\underline{k}}
\newcommand{\sheafn}{\underline{n}}
\newcommand{\sheafGamma}{\underline{\Gamma}}

\sisi{
\DeclareMathOperator{\simarrow}{\rule{.15mm}{1.8mm}\kern -3.4pt\rightsquigarrow}
\newcommand{\sheafell}{\underline{\ell}}
\newcommand{\sheafZZ}{\underline{\ZZ}}
}{
\newcommand{\nderef}[1]{la N.D.E.~(\ref{#1}) page~\pageref{#1}}
}

%%%%%%%%%%%%%%%%%
%%%%% expose XVII
%%%%%%%%%%%%%%%%%



\externaldocument{../01/01.tex}
\externaldocument{../02/02.tex}
\externaldocument{../03/03.tex}
\externaldocument{../04/04.tex}
\externaldocument{../05/05.tex}
\externaldocument{../05bis/05bis.tex}
\externaldocument{../06/06.tex}
\externaldocument{../07/07.tex}
\externaldocument{../08/08.tex}
\externaldocument{../09/09.tex}
\externaldocument{../10/10.tex}
\externaldocument{../11/11.tex}
%\externaldocument{../12/12.tex}
\externaldocument{../13/13.tex}
\externaldocument{../14/14.tex}
\externaldocument{../15/15.tex}
\externaldocument{../16/16.tex}
\externaldocument{../17/17.tex}
\externaldocument{../18/18.tex}
\externaldocument{../19/19.tex}

\setboolean{original}{false}

\begin{document}

\begin{center}
XII. Théorème de changement de base pour un morphisme propre\\
M. Artin 
\end{center}

\selectlanguage{english}version : \showgitstatus\selectlanguage{french}
\tableofcontents

\else
\chapter{Théorème de changement de base pour un morphisme propre}
%\addtocontents{toc}{. }
\begin{center}
M. Artin 
\end{center}
\fi

\setcounter{pageoriginal}{78}

\section{Introduction}\etiquette[1]{XII.sec1}\pageoriginale

Nous donnons dans cet exposé le théorème fondamental qui dit, sous une
forme un peu plus précise, que pour un morphisme \emph{propre}
$f:X\rightarrow Y$ de schémas, et pour un faisceau abélien $F$ de
torsion sur $X$, la fibre de $R^qf_{\ast}F$ en un point géométrique de
$Y$ est isomorphe à la cohomologie $H^q(X_{\xi}, F| X_{\xi})$ de la
fibre $X_{\xi}$.

Rappelons que pour un morphisme propre d'espaces paracompacts le
résultat analogue est bien connu (GODEMENT II 4.11) et facile. Il est
beaucoup plus délicat dans le cas des schémas. Par exemple,
l'hypothèse que $F$ soit de torsion est bien nécessaire, contrairement
à ce qui se passe pour les espaces paracompacts.

La démonstration se fait en plusieurs étapes: on se réduit par des
méthodes de dévissage plus ou moins formelles au cas où $Y$ est
noethérien, où la dimension relative de $f$ est $\leq 1$, et où le
faisceau $F$ est constant. Pour traiter ce cas particulier, on se sert
du théorème de spécialisation pour le groupe fondamental
(\hyperref[XII.thm5.9]{5.9}) qui est 
une variante non-abélienne du théorème énoncé plus haut, et qu'on peut
démontrer directement dans le cas particulier envisagé (\hyperrefext[XIII][XIII.sec2]{2}).

\section{Un exemple}\etiquette[2]{XII.sec2}\pageoriginale

Montrons que l'hypothèse que $F$ soit de torsion est nécessaire: Soit
$X$ une surface non singulière sur un corps $k$ algébriquement clos,
$Y$ une courbe non singulière sur $k$, et $f:X\rightarrow Y$ un
morphisme propre dont la fibre générique est lisse et la fibre
géométrique en le point géométrique $\bar{y}$ de $Y$ a un point double
ordinaire, et est irréductible. Soit $F$ le faisceau constant
$\ZZ_X$.

Rappelons le fait suivant: pour un point quelconque $z$ d'un schéma
$Z$, soit $i:z\rightarrow Z$ l'inclusion. On a
$H^1(Z,i_{\ast}\ZZ_z)=0$ (\hyperrefext[IX][IX.prop3.6]{3.6} (i)). Or puisque $X$ est
normal donc unibranche, il s'ensuit que $\ZZ_X\simeq
i_{\ast}\ZZ_x$, où $i:x\rightarrow X$ est l'inclusion du point
générique. Donc $H^1(X,\ZZ_X)=0$ et de même
$R^1f_{\ast}\ZZ_X=0$. Mais je dis que pour la fibre géométrique
$X_{\circ}$ de $X/Y$ avec le point double ordinaire on a
$H^1(X_{\circ},\ZZ_{X_{\circ}})\simeq \ZZ$\footnote[1]{Comparer, du point de vue $\Pi_1$, avec \SGA 3 X (6, Exemples) et (1.6).}, ce qui donne le contre-exemple, puisque
$\ZZ_{X_{\circ}}\simeq \ZZ_X| X_{\circ}$.

Soit $i_{\circ}:x_{\circ}\rightarrow X_{\circ}$ l'inclusion du point
générique de $X_{\circ}$, et soit $Q$ le point singulier de
$X_{\circ}$. La courbe $X_{\circ}$ a deux « branches »  en $Q$ et
il s'ensuit que la fibre de $i_{\circ^{\ast}}(\ZZ_{x_{\circ}})$
au point $Q$ est isomorphe à $\ZZ\times \ZZ$. En dehors
de $Q$ la courbe $X_{\circ}$ est normale et donc
$i_{\circ^{\ast}}(\ZZ_{x_{\circ}})\simeq
\ZZ_{X_{\circ}}$ dans $X_{\circ}-Q$. On a donc une suite exacte
$$
0\longrightarrow \ZZ_{X_{\circ}}\longrightarrow
i_{\circ^{\ast}}(\ZZ_{x_{\circ}})\longrightarrow
(\ZZ)_Q\longrightarrow 0{\quoi},
$$
où le dernier membre est l'extension par $0$ du faisceau $\ZZ$ au
point $Q$. Tous les trois membres ont un $H^{\circ}$ isomorphe à
$\ZZ$, et $H^1(X_{\circ},
i_{\circ^{\ast}}(\ZZ_{x_{\circ}}))=0$, d'où
$H^1(X_{\circ},\ZZ_{X_{\circ}})=\ZZ$\pageoriginale 
par la suite exacte de cohomologie.

\section{Rappels sur le $H^1$ non-abélien}\etiquette[3]{XII.sec3}

On va indiquer brièvement quelques résultats sur la cohomologie
non-abélienne dont nous aurons besoin dans la suite. Nous omettrons la
plupart des démonstrations, bien connues dans divers cas particuliers,
et qu'on trouvera dans la thèse de \textsc{Giraud} [1].

Soient $X$ un schéma, et $F$ un faisceau en groupes sur $X$. Pour
chaque $X'/X$ étale, on définit un ensemble pointé $H^1(X',F)$,
fonctoriel covariant en $F$ et contravariant en $X'$. Le $H^1$ commute
aux produits finis de faisceaux. Si $f:X\rightarrow Y$ est un
morphisme, on définit un faisceau d'ensembles pointés (c'est-à-dire,
avec section distinguée) $R^1f_{\ast}F$. C'est le faisceau associé au
préfaisceau (d'ensembles pointés) $\Rr^1F$, qui, à chaque
$Y'/Y$ étale, associe l'ensemble $H^1(\fprod{X}{Y'}{Y},F)$.

On a les suites exactes de cohomologie habituelles, dont nous nous
contenterons d'indiquer une:

\begin{enonce*}{Proposition 3.1}\etiquette[3.1]{XII.prop3.1}
Soit $0\rightarrow F\xrightarrow{u} G$ une injection de faisceaux de
groupes, et soit $C=G/F$ le faisceau d'ensembles homogènes sous
$G$. On a une suite exacte de faisceaux d'ensembles pointés
$$
0\rightarrow f_{\ast}F\rightarrow f_{\ast}G\rightarrow
f_{\ast}C\xrightarrow{\partial} R^1f_{\ast}F\xrightarrow{u^1}
R^1f_{\ast}G{\quoi}.
$$
\end{enonce*}

Bien entendu, cette suite exacte est obtenue en passant à la suite de
faisceaux associée à la suite exacte de préfaisceaux dont la valeur
pour $Y'/Y$ étale est
$$
0\rightarrow F(X')\rightarrow G(X')\rightarrow C(X')
\xrightarrow{\partial} H^1(X',F)\rightarrow H^1(X',G){\quoi},
$$
où\pageoriginale $X'=\fprod{X}{Y'}{Y}$. La relation d'équivalence
induite par $\partial$ est celle donnée par l'opération de $f_{\ast}G$
sur $F_{\ast}C$. On peut comme d'habitude expliciter la relation
d'équivalence induite par $u^1$ grâce à la notion de « torsion » 
[1].

\begin{enonce*}{Proposition 3.2}\etiquette[3.2]{XII.prop3.2}
Soient $X\xrightarrow{f}Y\xrightarrow{g}Z$ des morphismes de
schémas. On a une suite exacte pour chaque $Z'/Z$ étale, posant
$Y'=\fprod{Y}{Z'}{Z}$, $X'=\fprod{X}{Z'}{Z}$:
$$
0\longrightarrow H^1(Y',f_{\ast}F)\longrightarrow
H^1(X',F)\longrightarrow H^{\circ}(Y',R^1f_{\ast}F){\quoi},
$$
d'où en passant aux faisceaux associés, une suite exacte de faisceaux
pointés sur $Z$:
$$
0\longrightarrow (R^1g_{\ast})f_{\ast}F\longrightarrow
R^1(gf)_{\ast}F\longrightarrow g_{\ast}(R^1f_{\ast})F{\quoi}.
$$
\end{enonce*}

\begin{enonce*}{Proposition 3.3}\etiquette[3.3]{XII.prop3.3}
La famille des foncteurs de la forme $H^1(X',F)$, pour $X'$ étale sur
$X$, et de la forme $R^1f_{\ast}F$, pour un morphisme quelconque
$f:X\rightarrow Y$, est effaçable, c'est-à-dire il existe une
injection $F\xrightarrow{u}G$ tel que les morphismes induits
$H^1(X',F)\rightarrow H^1(X',G)$ et $R^1f_{\ast}F\rightarrow
R^1f_{\ast}G$ soient tous nuls. Plus précisément, il existe une
injection $F\xrightarrow{u}G$ tel que chacun des ensembles
$H^1(X',G)$, avec $X'$ étale sur $X$, et chacun des faisceaux pointés
$R^1f_{\ast}G$, soient nuls. Si $X$ est noethérien et $F$ est un
faisceau de ind-$\LL$-groupes constructible, alors il existe
une telle injection, où de plus $G$ est un faisceau de
ind-$\LL$-groupes.
\end{enonce*}

Démonstration. La première assertion est sans doute vraie dans un
topos quelconque, mais pour la topologie étale, il est commode de
faire l'effacement avec la « résolution de Godement » 
\begin{equation*}
F\longrightarrow G=\prod_{x\in X}i_{x^{\ast}}i_{x}{}^{\ast}F{\quoi},\tag{3.4}
\end{equation*}\etiquette[3.4]{XII.eq3.4}
où\pageoriginale $x$ parcourt l'ensemble des points de $X$ et où
$i_{x}:\bar{x}\rightarrow X$ est un point géométrique de $X$ au-dessus
de $x$. Notons que pour $X'/X$ étale, on a
$$
G(X')=\prod_{x\in
  X}(i_{x}{}^{\ast}F)(\fprod{X'}{\bar{x}}{X})=\prod_{x\in
  X}(\prod_{\bar{x}' /\bar{x}}F_{\bar{x}'}){\quoi},
$$
où les $\bar{x}'$ sont des points géométriques de $X'$. Or il est
clair que pour chaque point de $X'$ il y a au moins un $\bar{x}'$
au-dessus, d'où $F(X')\subset G(X')$ (VIII 3), donc $F\subset
G$. D'autre part, un $\bar{u}/\bar{x}$ étale quelconque est somme de
spectres de corps séparablement clos, d'où $H^1(\bar{u},\bar{G})=0$,
quel que soit le faisceau $\bar{G}$ sur $\bar{x}$. Il résulte alors de
\hyperref[XII.prop3.2]{3.2} que $H^1(X', i_{x^{\ast}}i_x{}^{\ast}F)=0$ pour chaque $X'/X$
étale, d'où (puisque $H^1$ commute aux produits) $H^1(X',G)=0$, donc
de même $R^1f_{\ast}G=0$ pour tout $f:X\rightarrow Y$, grâce à la
description de $R^1f_{\ast}(G)$ comme faisceau associé au préfaisceau
$Y'\mapsto H^1(\fprod{X}{Y'}{Y},G)$.

Supposons maintenant que de plus $F$ est constructible et $X$
noethérien. Alors on voit immédiatement qu'il existe un ensemble fini
de points $\{x_1,\ldots,x_n\}$ de $X$ tel que le morphisme
$$
F\longrightarrow \prod_{1\leq i\leq n}i_{x^{\ast}}i_{x}{}^{\ast}F
$$
soit déjà injectif. Le faisceau $i_{x^{\ast}}i_{x}{}^{\ast}F$ est
un faisceau de groupes ind-fini (IX 1.5), donc le produit fini l'est
aussi (\hyperrefext[IX][IX.prop1.4]{1.4} et \hyperrefext[IX][IX.prop1.6]{1.6} (iii)). 
En remplaçant $G$ par ce faisceau, on
obtient la dernière assertion.

\setcounter{subsection}{4}
\subsection{}\etiquette[3.5]{XII.subsec3.5}
Comme dans le cas de la cohomologie abélienne (traité dans V 2.4), on
peut calculer le $H^1$ par le procédé de \v{C}ECH [1].

\section{Le morphisme de changement de
  base}\etiquette[4]{XII.sec4}\pageoriginale 

Soit
$$
\xymatrix@C=2.75cm@R=1.5cm{
X   \ar[d]_{f}& X'\ar[l]_{g'}\ar[d]^{f'}\\
S       & S'\ar[l]
}
$$
un carré cartésien de morphisme de schémas. On en déduit un morphisme
de foncteur $\varphi=\varphi_g$
\begin{equation*}
g^{\ast}f_{\ast}\xrightarrow{\varphi}f'_{\ast}{g'}^{\ast}\tag{4.1}
\end{equation*}\etiquette[4.1]{XII.eq4.1}
de la manière suivante: donner un tel morphisme équivaut à donner un
morphisme de foncteurs
\begin{equation*}
f_{\ast}\longrightarrow g_{\ast}{f'}_{\ast}{g'}^{\ast}{\quoi},\tag{4.1 adj}
\end{equation*}\label{4.1 adj} 
parce que $g_{\ast}$ et $g^{\ast}$ sont adjoints. Nous prenons pour ce
morphisme le morphisme
$$
f_{\ast}\xrightarrow{(\id\rightarrow {g'}_{\ast}{g'}^{\ast})}
f_{\ast}{g'}_{\ast}{g'}^{\ast}=g_{\ast}{f'}_{\ast}{g'}^{\ast}{\quoi},
$$
en tenant compte que
$$
g_{\ast}{f'}_{\ast}=(gf')_{\ast}=(fg')_{\ast}=f_{\ast}{g'}_{\ast}{\quoi}.
$$

Notons qu'il y a un autre candidat pour \eqhyperref[XII.eq4.1]{4.1}. En effet, donner un tel
morphisme équivaut à donner
$$
{f'}^{\ast}g^{\ast}f_{\ast}\longrightarrow {g'}^{\ast}
$$
et on peut prendre
$$
{f'}^{\ast}g^{\ast}f_{\ast}\simeq
{g'}^{\ast}f^{\ast}f_{\ast}\xrightarrow{(f^{\ast}f_{\ast}\rightarrow
  \id)} {g'}^{\ast}{\quoi}.
$$

J'ignore\pageoriginale si ces deux choix coïncident toujours
\footnote[1]{Mr. \textsc{Deligne} a résolu par l'affirmative cette perplexité,
dans un contexte sensiblement plus général, \cf Exp. XVII.}, et en
tout cas nous nous servirons uniquement du premier. Nous l'appelons
\emph{morphisme de changement de base.}

En restreignant \eqhyperref[XII.eq4.1]{4.1} aux faisceaux abéliens (\resp de groupes) on
déduit un morphisme $\varphi^q=\varphi^g g$ pour chaque $q\geq 0$
(\resp pour $q=0$, 1):
\begin{equation*}
g^{\ast}(R^qf_{\ast})\xrightarrow{\varphi^q}(R^q{f'}_{\ast}){g'}^{\ast}{\quoi}.
\tag{4.2}
\end{equation*}\etiquette[4.2]{XII.eq4.2}
En effet, cela revient au même que de donner
\begin{equation*}
(R^qf_{\ast})\longrightarrow
  g_{\ast}(R^q{f'}_{\ast}){g'}^{\ast}{\quoi},\tag{4.2 adj}
\end{equation*}\label{4.2 adj}
et on a des morphismes « évidents » 
\begin{gather*}
(R^qf_{\ast})\xrightarrow{a}
  (R^qf_{\ast}){g'}_{\ast}{g'}^{\ast}\xrightarrow{b}
  R^q{'}fg')_{\ast}{g'}^{\ast}=\\
=(R^q(gf')_{\ast}{g'}^{\ast}\xrightarrow{c}
  g^{\ast}(R^q{f'}_{\ast}){g'}^{\ast}{\quoi}.
\end{gather*}

\begin{enonce*}[remark]{Remarque 4.3}\etiquette[4.3]{XII.rem4.3}
Supposons que $g$ (donc $g'$) soit une immersion fermée. Alors les
morphismes $b$ et $c$ ci-dessus sont des isomorphismes (VIII 5.7). Il
suit que \eqhyperref[XII.eq4.2]{4.2} n'est autre essentiellement que le morphisme a
ci-dessus, qui est défini en appliquant $R^qf_{\ast}$ au morphisme
canonique $F\rightarrow {g'}_{\ast}{g'}^{\ast}F$.
\end{enonce*}

\begin{enonce*}{Proposition 4.4}\etiquette[4.4]{XII.prop4.4}
(Composition des morphismes de changement de base).
\begin{enumerate}
\renewcommand{\theenumi}{\roman{enumi}}
\renewcommand{\labelenumi}{(\theenumi)}
\item Soit
$$
\xymatrix@C=2.5cm@R=1.5cm{
X   \ar[d]_{f}& X' \ar[l]_{g'} \ar[d]_{f'} & X'' \ar[l]_{h'}\ar[d]_{f''}\\
S   & S' \ar[l]_{g} &  S''\ar[l]_{h}
}
$$
un\pageoriginale diagramme cartésien de morphismes de schémas. Alors
le diagramme suivant est commutatif:
$$
\xymatrix@C=2cm@R=1.5cm{
(gh)^{\ast}f_{\ast}  \ar[rr]^{gh} \ar[d]^{\text{\rotatebox{90}{$\backsim$}}}& &f''_{\ast}(g'h')^{\ast}\ar[d]^{\text{\rotatebox{90}{$\backsim$}}}\\
h^{\ast}g^{\ast}f_{\ast} \ar[r]^{h^{\ast}{}_{\ast}g} &
h^{\ast}f'_{\ast}g'^{\ast} \ar[r]^{h ^{}\ast g'^{\ast}} &
f''_{\ast}h'^{\ast}g'^{\ast}}
$$
et le diagramme avec $f_{\ast}$, $f'_{\ast}$, $f''_{\ast}$ remplacé par
$R^qf_{\ast}$, $R^q{f'}_{\ast}$, $R^q{f''}_{\ast}$ l'est aussi, dans
les domaines de définition naturels des foncteurs en question {\rm
  (\cf ${\rm n^{\circ}}$ 3).}

\item Soit
$$
\xymatrix@C=3.25cm@R=1.25cm{
Y \ar[d]_{e}& Y'\ar[l]_{g''} \ar[d]^{e'}\\
X \ar[d]_{f}& X'\ar[l]_{g'}\ar[d]^{f'}\\
S & S'\ar[l]_{g}
}$$
un diagramme cartésien de morphismes de schémas. Alors le diagramme
suivant est commutatif:
$$
\xymatrix@C=2.75cm@R=1.5cm{
g^{\ast}(fe)_{\ast} \ar[rr]\ar@{=}[d]& & (f'e')_{\ast}g''^{\ast}\ar@{=}[d]\\
g^{\ast}f_{\ast}e_{\ast}   \ar[r] &f'_{\ast}g'{}^{\ast}e_{\ast} \ar[r] 
&f'_{\ast}e'_{\ast}g''^{\ast},
}
$$
où les trois flèches horizontales proviennent des homomorphismes de
changement de base pour $fe$, $f$ et $e$ respectivement.

Dans le cas abélien, les morphismes de changement de base pour $fe$
sont induits par un morphisme de suites spectrales
$$
\xymatrix{
E^{p,q}_2 = g^{\ast}(R^p f_{\ast})(R^{q} e_{\ast})F \ar@{=>}[r] \ar[d]_{\varphi}&
g^{\ast}R^{p+q}(fe)_{\ast}F\ar[d]_{{\varphi}}\\
(R^{p}f'_{\ast})(R^q e'_{\ast})g''^{\ast}F \ar@{=>}[r]& R^{p+q}(f'e')_{\ast}g''^{\ast}F
}
$$
qui,\pageoriginale pour les termes initiaux, est induit par les
morphismes de changement de base pour $e$ et $f$; et pour un faisceau
de groupes, on a un morphisme de suites exactes
$$
\xymatrix{
0  \ar[r] & g^{\ast}(R^{1}f_{\ast})e_{\ast}F \ar[r]\ar[d]^{\varphi} &
g^{\ast}R^{1}(fe)_{\ast}F \ar[r] \ar[d]^{\varphi}&
g^{\ast}f_{\ast}(R^{1}e_{\ast})F \ar[d]^{\varphi}\\
0 \ar[r]& ((R^{1}f'_{\ast})e'_{\ast}g''^{\ast}F) \ar[r] &
R^1(f' e')_{\ast}g''^{\ast}F  \ar[r] & f'_{\ast}(R^{1}e'_{\ast})g''^{\ast}F
}
$$
\end{enumerate}
\end{enonce*}

Nous laissons la démonstration au lecteur.

\section{Énoncé du théorème principal et de quelques variantes}\etiquette[5]{XII.sec5}

Le théorème est le suivant:

\begin{enonce*}{Théorème 5.1}\etiquette[5.1]{XII.thm5.1}
(Théorème de changement de base pour un morphisme propre). Soit
$$
\xymatrix@C=3.25cm@R=1.5cm{
X \ar[d]_{f} &   X' \ar[l]_{g'} \ar[d]^{f'}\\
S                      & S' \ar[l]_{g}
}
$$
un diagramme cartésien de schémas, avec $f$ propre.
\begin{enumerate}
\renewcommand{\theenumi}{\roman{enumi}}
\renewcommand{\labelenumi}{(\theenumi)}
\item Soit $F$ un faisceau d'ensembles sur $X$. Alors le morphisme
  \eqhyperref[XII.eq4.1]{4.1}
$$
\varphi:g^{\ast}f_{\ast}F\longrightarrow {f'}_{\ast}{g'}^{\ast}F
$$
est bijectif.

\item Soit\pageoriginale $f$ un faisceau de groupes 
(\resp de groupes ind-finis) sur
  $X$. Alors le morphisme \eqhyperref[XII.eq4.2]{4.2}
$$
\varphi^1:g^{\ast}(R^1f_{\ast})F\rightarrow
(R^1{f'}_{\ast}){g'}^{\ast}F
$$
est injectif (\resp bijectif).

\item Soit $F$ un faisceau abélien de torsion sur $X$. Alors le
  morphisme \eqhyperref[XII.eq4.2]{4.2}
$$
\varphi^q:g^{\ast}(R^qf_{\ast})F\longrightarrow
(R^q{f'}_{\ast}){g'}^{\ast}F
$$
est bijectif pour chaque $q\geq 0$.
\end{enumerate}
\end{enonce*}

On va d'abord énoncer des conséquences immédiates du théorème. Si
$S'=\xi$ est un point géométrique de $S$, de sorte qu'un faisceau sur
$\xi$ est déterminé par ses sections globales, on a
$H^{\circ}(\xi,R^q{f'}_{\ast}G)=H^q(X',G)$ pour chaque faisceau $G$
sur $X$, d'où le corollaire suivant (qui est d'ailleurs
essentiellement équivalent à (\hyperref[XII.thm5.1]{5.1})):

\begin{enonce*}{Corollaire 5.2}\etiquette[5.2]{XII.cor5.2}
Soient $f:X\rightarrow S$ propre, $\xi\rightarrow S$ un point
géométrique, et $X_{\xi}$ la fibre de $f$ au point $\xi$.
\begin{enumerate}
\renewcommand{\theenumi}{\roman{enumi}}
\renewcommand{\labelenumi}{(\theenumi)}
\item Si $F$ est un faisceau d'ensembles sur $X$, le morphisme
  canonique de changement de base
$$
(f_{\ast}F)_{\xi}\longrightarrow H^{\circ}(X_{\xi},F| X_{\xi})
$$
est bijectif.

\item Si $F$ est un faisceau de groupes (\resp de groupes ind-finis)
  le morphisme de changement de base
$$
R^1f_{\ast}(F)_{\xi}\longrightarrow H^1(X_{\xi},F| X_{\xi})
$$
est injectif (\resp bijectif).

\item Si\pageoriginale $F$ est un faisceau abélien de torsion, alors
  pour chaque $q\geq 0$ le morphisme de changement de base
$$
(R^qf_{\ast})F_{\xi}\longrightarrow H^q(X_{\xi},F| X_{\xi})
$$
est bijectif.
\end{enumerate}
\end{enonce*}

Tenant compte de (\hyperrefext[X][X.cor4.3]{4.3}) et (\hyperrefext[X][X.cor5.2]{5.2}), on trouve:

\begin{enonce*}{Corollaire 5.3}\etiquette[5.3]{XII.cor5.3}
Sous les conditions de (\hyperrefext[X][X.cor5.2]{5.2}), soit $n$ la dimension de la fibre
$X_{\xi}$. Alors pour chaque faisceau abélien de torsion $F$ sur $X$,
la fibre $(R^qf_{\ast}F)_{\xi}$ est nulle pour $q>2n$. Si $F$ est un
faisceau de $p$-torsion ($p$ la caractéristique de $\xi$), alors la
fibre de $R^qf_{\ast}(F)$ en $\xi$ est nulle pour $q>n$.
\end{enonce*}

Rappelons qu'on dit que la \emph{dimension relative de $X/S$ est}
$\leq n$ si la dimension de chaque fibre est $\leq n$. On a donc

\begin{enonce*}{Corollaire 5.3 bis}\etiquette[5.3 bis]{XII.corbis5.3}
Soit $f:X\rightarrow S$ propre et $X/S$ de dimension relative $\leq
n$. Alors $R^qf_{\ast}F=0$ si $q>2n$ pour chaque faisceau abélien $F$
de torsion sur $X$. Si $S$ est de caractéristique $p$, et si $F$ est
un faisceau abélien de $p$-torsion, alors $R^qf_{\ast}F=0$ si $q>n$.
\end{enonce*}

En appliquant (\hyperref[XII.thm5.1]{5.1}) a un morphisme de spectres de corps séparablement
clos, on obtient

\begin{enonce*}{Corollaire 5.4}\etiquette[5.4]{XII.cor5.4}
(Invariance de la cohomologie par changement de corps de base dans le
  cas propre). Soient $k\subset K$ des corps séparablement clos et $X$
  un schéma propre sur $\Spec k$. Soit $X'=X_K$. Soit $F$ un faisceau
  sur $X$ et $F'$ le faisceau image inverse sur $X'$. Alors
\begin{enumerate}
\renewcommand{\theenumi}{\roman{enumi}}
\renewcommand{\labelenumi}{(\theenumi)}
\item $H^{\circ}(X,F)\xrightarrow{\sim} H^{\circ}(X',F')$.

\item Si\pageoriginale $F$ est un faisceau de groupes (\resp de
  groupes ind-finis)
$$
H^1(X,F)\xrightarrow{\sim} H^1(X',F')
$$
est injectif (\resp bijectif).

\item Si $F$ est un faisceau abélien de torsion,
$$
H^q(X,F)\xrightarrow{\sim} H^q(X',F')
$$
pour tout $q$.
\end{enumerate}
\end{enonce*}

Une variante légèrement différente de (\hyperref[XII.cor5.2]{5.2}) est la suivante:

\begin{enonce*}{Corollaire 5.5}\etiquette[5.5]{XII.cor5.5}
Soit $S$ le spectre d'un anneau hensélien et soit $S'=s_{\circ}$ le
point fermé. Soit $f:X\rightarrow S$ propre. Posons $X_{\circ}=X'$,
$F_{\circ}={g'}^{\ast}F$. Alors on a:
\begin{enumerate}
\renewcommand{\theenumi}{\roman{enumi}}
\renewcommand{\labelenumi}{(\theenumi)}
\item $H^{\circ}(X,F)\xrightarrow{\sim}
  H^{\circ}(X_{\circ},F_{\circ})$

pour tout faisceau d'ensembles $F$.

\item $H^1(X,F)\rightarrow H^1(X_{\circ},F_{\circ})$

est injectif (\resp bijectif) pour tout faisceau de groupes (\resp de
groupes ind-finis)~$F$.

\item $H^q(X,F)\xrightarrow{\sim} H^q(X_{\circ},F_{\circ})$

pour tout faisceau abélien de torsion $F$, et tout $q$.
\end{enumerate}
\end{enonce*}

En effet, cet énoncé n'est qu'un cas particulier de (\hyperref[XII.cor5.2]{5.2}) si $S$ est
strictement local. Supposons que (\hyperref[XII.cor5.2]{5.2}) soit connu. Soit $\bar{S}$ le
localisé strict de $S$, $\bar{s}$ le point fermé de $\bar{S}$, et $G$
le groupe de Galois de $\bar{S}/S$ (qui est aussi le groupe de Galois
de $\bar{s}/s_{\circ}$). Alors $G$ opère sur les fibres des $R^q$, et
on a (avec des notations évidentes)
\begin{align*}
H^{\circ}(X,F) &\simeq H^{\circ}(\bar{X},\bar{F})^G\\
H^{\circ}(X_{\circ},F_{\circ}) &\simeq H^{\circ}(\bar{X}_{\circ}, \bar{F}_{\circ})^G
\end{align*}
pour\pageoriginale un faisceau d'ensembles $F$, d'où (i) puisque
$H^{\circ}(\bar{X},\bar{F})\xrightarrow{\sim}H^{\circ}(\bar{X}_{\circ},\bar{F}_{\circ})$et que l'isomorphisme commute évidemment avec l'opération de $G$.

Si $F$ est un faisceau de groupes, on définit aisément un morphisme de
suites exactes d'ensembles pointés
$$\vcenter{\etiquette[5.5.1]{XII.eq5.5.1}
\xymatrix{
0  \ar[r] & H^{1}(G,H^{0}(\bar{X},\bar{F})) \ar[r]\ar[d]& H^{1}(X,F) \ar[r]\ar[d] &
H^{1}(\bar{X},\bar{F})^G \ar[d]\\
0  \ar[r] & H^{1}(G,H^{0}(\bar{X}_0,\bar{F}_0)) \ar[r]& H^{1}(X_0,F_0) \ar[r] &
H^{1}(\bar{X}_0,\bar{F}_0)^G ,
}}\leqno{(5.5.1)}
$$
où les flèches verticales extrêmes sont des isomorphismes d'après
(\hyperref[XII.cor5.2]{5.2}). Soient $\alpha$, $\beta\in H^1(X,F)$ des classes
qui ont même 
image dans $H^{\circ}(X_{\circ},F_{\circ})$, et soit $F^{\alpha}$ le
faisceau obtenu en tordant $F$ à l'aide de $\alpha$ \cite{SGA4-12:1}. Soit $\beta'$
la classe $H^1(X,F^{\alpha})$ correspondant à $\beta$. Alors, puisque
la classe $\alpha'$ correspondant à $\alpha$ est nulle, l'image de
$\beta'$ dans $H^1(X_{\circ},(F^{\alpha})_{\circ})$ est nulle.

Examinons le diagramme déduit de \eqhyperref[XII.eq5.5.1]{5.5.1} en remplaçant $F$ par
$F^{\alpha}$: on trouve que l'image de $\beta'$ dans
$H^{\circ}(\bar{X},\bar{F}^{\alpha})^G$ est nulle, donc que $\beta'$
est image d'un élément de $H^1(G,H^{\circ}(\bar{X},\bar{F}))$, et que
cet élément doit être nul. Donc $\beta'$ est nul, \ie
$\beta'=\alpha'$, d'où $\beta=\alpha$. Cela donne l'injectivité du
morphisme de (ii).

Supposons maintenant que $F$ soit un faisceau de groupes ind-finis, et
démontrons la surjectivité de la flèche (ii). Or il suffit de
démontrer que le foncteur $F\mapsto H^1(X_{\circ},F_{\circ})$ est
effaçable. En effet, soit $0\rightarrow F\xrightarrow{u}G$ une
injection qui efface ce foncteur. Soit $C=G/F$. On a un morphisme de
suites exactes (\hyperref[XII.prop4.4]{4.4})
$$
\xymatrix{
H^{0}(X,C)      \ar[r]\ar[d] & H^{1}(X,F) \ar[d]_{\epsilon}         \ar[r]         &
H^{1}(X,G) \ar[d]\\
H^{0}(X_0,C_0) \ar[r] & H^{1}(X_{0},F_{0})  \ar[r]^{u_0^1} & H^{1}(X_0,G_0)
}
$$
où le morphisme $u^1_{\circ}$ est nul, ce qui implique bien la
surjectivité de $\epsilon$.

Démontrons\pageoriginale que le foncteur $F\mapsto
H^1(X_{\circ},F_{\circ})$ est effaçable pour les faisceaux $F$ de
groupes ind-finis. Soit $F\xrightarrow{u}C$ la « résolution de
Godement »  \eqhyperref[XII.eq3.4]{3.4}. On a
$$
H^1(G,H^{\circ}(\bar{X}_{\circ},
\bar{C}_{\circ}))\simeq H^1(G,H^{\circ}(\bar{X},\bar{C}))\subset
H^1(X,C)=0{\quoi}.
$$

De plus, on a le carré commutatif:
$$
\xymatrix{
H^{1}(\bar{X},\bar{F}) \ar@{-}[r]^{\sim} \ar[d]& H^{1}(\bar{X}_0,\bar{F}_0)\ar[d]\\
0 = H^{1}(\bar{X},\bar{C})  \ar[r]     & H^{1}(\bar{X}_0,\bar{C}_0)
}
$$
(puisque $H^1(\bar{X},\bar{G})=\varinjlim H^1(X',C)=0$, $X'/X$ étale),
donc est nul. On a donc un morphisme de suites exactes
$$
\xymatrix{
0 \ar[r] & H^{1}(G,H^{0}(\bar{X}_0, \bar{F}_0))  \ar[r] \ar[d]&
H^{1}(X_0,F_0) \ar[r] \ar[d]^{\beta}& H^{1}(\bar{X}_0,\bar{F}_0)^G \ar[d]^{\alpha^G}\\
&0  \ar[r] & H^{1}(X_0,G_0)  \ar[r] & H^{1}(\bar{X}_0,\bar{G}_0)^G
}
$$
où $\alpha^G$ est nul, ce qui implique que $\beta$ est également nul,
d'où l'effaçabilité.

Démontrons enfin l'assertion (iii) de (\hyperref[XII.cor5.5]{5.5}): On a un morphisme de
suites spectrales
$$
\xymatrix{
**[l]E_2^{p,q} = H^{p}(G,H^{q} (\bar{X},\bar{F})  \ar@{=>}[r] \ar[d]& H^{n}(X,F) \ar[d]^{\varphi}\\
H^{p}(G,H^q,\bar{X}_0,\bar{F}_0 ))     \ar@{=>}[r] &  H^{n}(\bar{X}_0,\bar{F}_0)&.
}
$$
D'après (\hyperref[XII.cor5.2]{5.2}), les morphismes $H^q(\bar{X},\bar{F})\rightarrow
H^q(\bar{X}_{\circ},\bar{F}_{\circ})$ sont des isomorphismes. Cela
implique que $\varphi$ est un isomorphisme pour $E^{p,q}_2$, donc
pour les aboutissements, d'où le résultat, \cqfd

\setcounter{subsection}{5}
\subsection{}\etiquette[5.6]{XII.subsec5.6}
Examinons\pageoriginale (\hyperref[XII.cor5.5]{5.5}) (i) et (ii) dans le cas où $F$ est
constant: Soit $F=S_X$ un faisceau d'ensembles constant sur le schéma
$X$, de valeur $T$. Alors $F(X)$ n'est autre que l'ensemble des
applications localement constantes de $X$ dans $T$. Si $X_{\circ}$ est
un sous-schéma de $X$, et si $T$ a au moins deux éléments, on en
conclut que $F(X)\rightarrow F(X_{\circ})$ est bijectif si et
seulement si l'application $U\mapsto U_{\circ}=U\cap X_{\circ}$ de
l'ensemble $\Theta(X)$ des parties à la fois ouvertes et fermées de
$X$ dans l'ensemble analogue $\Theta(X_{\circ})$ est bijective. Si
l'ensemble $\Pi_{\circ}(X_{\circ})$ des composantes connexes de
$X_{\circ}$ est fini, cela revient au même que de dire que
l'application $\Pi_{\circ}(X_{\circ})\rightarrow \Pi_{\circ}(X)$ est
bijective. Ainsi, (\hyperref[XII.cor5.5]{5.5}) (i) peut aussi se reformuler de la façon
suivante (compte tenu que dans ce cas $X_{\circ}$ est de type fini sur
un corps, donc est noethérien, donc $\Pi_{\circ}(X_{\circ})$ est
fini):

\begin{enonce*}{Corollaire 5.7}\etiquette[5.7]{XII.cor5.7}
Soient $S$ le spectre d'un anneau hensélien, $s_{\circ}$ son point
fermé. Soit $f:X\rightarrow S$ un morphisme propre et $X_{\circ}$ la
fibre fermée de $X/S$. Alors l'application $\Pi_{\circ}(X)\rightarrow
\Pi_{\circ}(X_{\circ})$ induite par $X_{\circ}\rightarrow X$ est
bijective.
\end{enonce*}

(En particulier, cela implique que l'ensemble $\Pi_{\circ}(X)$ est
fini). Nous allons démontrer directement (\hyperref[XII.cor5.7]{5.7}) dans le cas noethérien:

\begin{enonce*}{Lemme 5.8}\etiquette[5.8]{XII.lem5.8}
L'énoncé $(5.7)$ est vrai si $S$ est noethérien \footnote[1]{\Cf \EGA
  IV 18.5.19.}.
\end{enonce*}

Démonstration. Dans le cas noethérien, $\Pi_{\circ}(X)$ est
certainement fini, et en remplaçant $X$ par une composante connexe
on peut supposer que $\Pi_{\circ}(X)$ est un ensemble à un élément,
c'est-à-dire que $X$ est connexe et non-vide. Il faut démontrer
qu'alors $X_{\circ}$ est également connexe et non-vide. Or l'image de
$X$ dans $S$ est fermée puisque $f$ est propre et non-vide car $X$ est
non-vide, donc contient $S_{\circ}$, ce qui explique que $X_{\circ}$
est non-vide. Soit
$$
\xymatrix@R=1cm@C=1.5cm{
X \ar[rd]^{f'} \ar[dd]_{f}& \\
  & S'\ar[ld]^{\pi}\\
S & 
}
$$
la\pageoriginale factorisation de Stein de $f$ (\EGA III 4.3.3) de
sorte que $S'/S$ est fini et $S'$ est connexe. On sait que chaque
fibre de $f'$ est connexe non-vide, et l'espace sous-jacent à
$X_{\circ}$ est la réunion des espaces sous-jacents aux fibres fermées
de $S'$. Tout revient donc à démontrer que $S'$ est le spectre d'un
anneau \emph{local}. Mais puisque $S$ est hensélien, l'anneau $A'$ de
$S'$ est produit d'anneaux locaux (VII 4.1 (i)), donc $A'$ est local
puisque $S'$ est connexe, ce qui prouve \hyperref[XII.lem5.8]{5.8}.

\setcounter{subsection}{9}
\setcounter{subsubsection}{-1}
\subsubsection{}\etiquette[5.9.0]{XII.subsubsec5.9.0}
Supposons maintenant que $F=G_X$ est un faisceau constant de groupes,
de valeur $G$, où $G$ est un groupe fini. Alors $H^1(X,G_X)$ classifie
les revêtements principaux galoisiens de $X$, de groupe $G$ (VII
2.1)). Le corollaire \hyperref[XII.cor5.5]{5.5} (ii) dit que le morphisme
$X_{\circ}\rightarrow X$ induit une bijection entre les ensembles de
classes de ces revêtements. Il s'ensuit donc, de la définition des
$\Pi_1$, que \hyperref[XII.cor5.5]{5.5} (ii) pour les faisceaux de groupes finis constants
(en présence de (i)) équivaut à:

\begin{enonce*}{Théorème 5.9}\etiquette[5.9]{XII.thm5.9}
(Théorème de spécialisation pour le groupe fondamental). Avec les
  notations de $(5.5)$, supposons $X$ connexe non-vide. Alors
  $X_{\circ}$ est connexe non-vide, et si $\bar{x}_{\circ}$ est un
  point géométrique de $X_{\circ}$, le morphisme de groupes
  fondamentaux
$$
\Pi_1(X_{\circ},\bar{x}_{\circ})\longrightarrow
\Pi_1(X,\bar{x}_{\circ})
$$
est bijectif.
\end{enonce*}

D'autre\pageoriginale part, d'après la théorie de Galois (\SGA V)
l'énoncé \hyperref[XII.thm5.9]{5.9} (dans le cas $X_{\circ}$ connexe, ce qu'on peut supposer,
grâce à \hyperref[XII.cor5.7]{5.7}) équivaut à

\begin{enonce*}{Théorème 5.9 bis}\etiquette[5.9 bis]{XII.thmbis5.9}
Soient $S$ le spectre d'un anneau hensélien, $s_{\circ}$ le point
fermé de $S$, $f:X\rightarrow S$ un morphisme propre, et $X_{\circ}$
la fibre fermée de $X/S$. Pour un schéma $Z$, notons $\Et(Z)$ la
catégorie des revêtements étales de $Z$. Alors le foncteur restriction
$$
\Et(X)\longrightarrow \Et(X_{\circ})
$$
est une équivalence de catégories.
\end{enonce*}

\begin{enonce*}[remark]{Remarque 5.10}\etiquette[5.10]{XII.rem5.10}
Rappelons que \hyperref[XII.thm5.9]{5.9} a été déjà démontré (\SGA 1 X 2.1) dans le cas où
$S$ est le spectre d'un anneau local noethérien
complet. Malheureusement on n'arrive pas à utiliser ce résultat pour
prouver \hyperref[XII.thm5.9]{5.9} en général \footnote[1]{La situation a changé depuis ces
  lignes ont été écrites, \cf th.~\hyperref[XII.prop3.1]{3.1} dans M. Artin, Algebraic
  approximation of structures over complete local rings (à
  paraître).}, et la démonstration donnée de 5.9, dans le présent
exposé et le suivant, est très différente de celle donnée dans \SGA 1 X.
\end{enonce*}

\setcounter{subsection}{10}
\subsection{}\etiquette[5.11]{XII.subsec5.11}
On peut donner une autre variante « non commutative »  de (\hyperref[XII.cor5.5]{5.5})
(ii), en utilisant la théorie de la 2-cohomologie de J. \textsc{Giraud} (\cf
Thèse de J. \textsc{Giraud}, en préparation). Supposons $S$ noethérien. Prenons
un \emph{lien} $L$ sur $X$, et supposons que $L$ se réalise localement
(pour la topologie étale) par un faisceau en groupes
ind-finis. Considérons l'application
$$
\varphi^{(2)}:H^2(X,L)\longrightarrow H^2(X_{\circ},L_{\circ}){\quoi},
$$
où $L_{\circ}$ est la restriction de $L$ à $X_{\circ}$. Alors on a
\begin{equation*}
(\varphi^{(2)})^{-1}(H^2(X_{\circ},L_{\circ})')=H^2(X,L)'{\quoi},\tag{5.11.1}
\end{equation*}\etiquette[5.11.1]{XII.eq5.11.1}
où\pageoriginale les « primes »  désignent les parties formées des
éléments neutres.

Compte tenu de \hyperref[XII.cor5.5]{5.5} (iii), cet énoncé équivaut aussi au suivant:

\emph{Pour toute gerbe $\mathcal{G}$ sur $X$, de lien $L$ comme ci-dessus,
  désignant par $\mathcal{G}_{\circ}$ la restriction de $\mathcal{G}$ à
  $X_{\circ}$, le foncteur restriction sur les catégories de sections}
$$
\mathcal{G}(X)\longrightarrow \mathcal{G}_{\circ}(X_{\circ})
$$
\emph{est une équivalence de catégories.}

Le fait que ce foncteur soit pleinement fidèle résulte en effet de \hyperref[XII.cor5.5]{5.5}
(i), et le fait qu'il soit essentiellement surjectif lorsque
$\mathcal{G}(X)\neq ∅$ \ie lorsque la classe $\xi \in H^2(X,\sisi{\mathcal{G}}{L})$
de $\mathcal{G}$ est neutre, résulte de \hyperref[XII.cor5.5]{5.5} (ii). 
Il reste donc à exprimer que 
$\mathcal{G}_{\circ}(X_{\circ})\neq ∅$ implique $\mathcal{G}(X)\neq
∅$, ce  qui n'est autre que \eqhyperref[XII.eq5.11.1]{5.11.1}.

Nous ne donnerons pas ici la démonstration de \eqhyperref[XII.eq5.11.1]{5.11.1}, et nous nous
contenterons de signaler que pour tout schéma noethérien $X$ et tout
sous-schéma fermé, on peut montrer en fait que la validité des énoncés
\hyperref[XII.cor5.5]{5.5} (i) (ii) implique celle de \eqhyperref[XII.eq5.11.1]{5.11.1}; la démonstration se fait
assez élémentairement, en utilisant le résultat de descente VIII 9.4
a).

\section{Premières réductions}\etiquette[6]{XII.sec6}

Nous utilisons l'abréviation $(f,F,g)$ pour désigner les données de
\hyperref[XII.thm5.1]{5.1} dans l'un quelconque des trois cas (i), (ii) ou (iii) de ce
théorème.

\begin{enonce*}{Lemme 6.1}\etiquette[6.1]{XII.lem6.1}
Pour que $5.1$ soit vrai pour les données $(f,F,g)$, où $f$, $F$ sont
fixés, et $g$ quelconque, il suffit que pour tout $S$-schéma
$\bar{S}$, localisé strict d'un schéma localement\pageoriginale de
type fini $S'$ sur $S$ en un point qui est fermé dans sa fibre, le
théorème soit vrai pour $(\bar{f},\bar{F},\bar{g})$, où $\bar{f}$,
$\bar{F}$ sont déduits de $f$, $F$ par le changement de base
$\bar{S}\rightarrow S$, et où $\bar{g}:\bar{s}\rightarrow \bar{S}$ est
l'inclusion du point fermé de $\bar{S}$.
\end{enonce*}

\noindent
\emph{Démonstration.} On peut supposer évidemment $S$, $S'$ affines,
spectres d'anneaux $A$, $A'$. Alors $A'$ est la limite inductive de
ses sous-$A$-algèbres de type fini $A'_i$, et utilisant la théorie de
passage à la limite (VII 5.11 et 5.14), on est ramené à prouver \hyperref[XII.thm5.1]{5.1}
pour les $(f,F,g_i)$ où $g_i$: $\Spec(A'_i)=S'_i\rightarrow S$, ce qui
nous ramène au cas où $g$ est de type fini. Pour prouver que les
homomorphisme de changement de base sont des isomorphismes, il suffit
de prouver qu'ils induisent des isomorphismes sur les fibres des
faisceaux en jeu, et il suffit de regarder les fibres géométriques en
des points $s'$ de $S'$ qui sont fermés dans leur fibre sur $S$ (VIII
3.13). Soit $\bar{s}'$ un point géométrique de $S'$ correspondant à
une clôture séparable de $k(s')$; comme $s'$ est algébrique sur le
point correspondant $s$ de $S$, $k(\bar{s}')$ est donc aussi une
clôture séparable de $k(s)$. Désignons par $\bar{s}$ le préschéma
$\bar{s}'$, considéré comme étant au-dessus de $S$, \ie considéré
comme point géométrique de $S$, et considérons les localisés stricts
$\bar{S}$ et $\bar{S}'$ de $S$ en $\bar{s}$ et en $\bar{s}'$
respectivement, de sorte qu'on a un carré commutatif.
$$
\vcenter{
\xymatrix@C=3cm@R=1.25cm{
\bar{s}'  \ar[r] \ar[d]& \bar{s}\ar[d]\\
\bar{S}'  \ar[r] &  \bar{S},
}}
$$
où la première flèche horizontale est un isomorphisme. Compte tenu de
VIII 5.2 et 5.3, l'homomorphisme induit sur les fibres en $\bar{s}'$
par l'homomorphisme de changement de base relatif à $(f,F,g)$
s'identifie à l'homomorphisme 
$$
H^{i}(\bar{X},\bar{F})\longrightarrow H^{i}(\bar{X}',\bar{F}')
$$
induit\pageoriginale par $\bar{X}'\rightarrow \bar{X}$, où $\bar{X}$ 
et $\bar{X}'$
désignent les schémas déduits de $X$ par les changements de base
$\bar{S}\rightarrow S$ et $\bar{S}'\rightarrow S$. Or le carré
commutatif ci-dessus fournit de façon évidente un carré commutatif
$$
\xymatrix@C=3cm@R=1.5cm{
H^{i}(\bar{X}_{\bar{s}}, \bar{F}) \ar[r]& H^{i}(\bar{X}'_{\bar{s}'},\bar{F}')
\\
H^{i}(\bar{X},\bar{F})  \ar[r] \ar[u]& H^{i}(\bar{X}',\bar{F}')
\ar[u],
}
$$
où la première flèche horizontale est encore un isomorphisme, donc la
deuxième flèche horizontale est un isomorphisme dès que les deux
flèches verticales le sont. Or c'est ce qu'assure l'hypothèse de
(\hyperref[XII.lem6.1]{6.1}), \cqfd

\setcounter{subsection}{1}
\subsubsection{}\etiquette[6.1.1]{XII.subsubsec6.1.1}
Ainsi, pour prouver \hyperref[XII.thm5.1]{5.1} dans une situation $(F,f,-)$, on est ramené à
étudier le cas particulier \hyperref[XII.cor5.5]{5.5}, avec $S$ strictement local. Cela nous
amène à étudier de façon générale la situation où on a un morphisme
$h:Y\rightarrow X$ de schémas, et où on se propose de donner les
conditions générales moyennant lesquelles les homomorphismes
correspondants
$$
H^{i}(X,F)\longrightarrow H^{i}(Y,h^{\ast}(F))
$$
sont bijectifs. Le présent numéro donnera quelques résultats
auxiliaires faciles sur cette situation, que (pour la commodité de
références ultérieures) nous énoncerons avec plus de généralité et de
précision qu'il ne serait nécessaire pour la démonstration de \hyperref[XII.thm5.1]{5.1}.

\begin{enonce*}{Lemme 6.2}\etiquette[6.2]{XII.lem6.2}
Soient $C$, $C'$ deux catégories abéliennes, $T^{\bigdot}=(T^{i})$ et
$T'^{\bigdot}=({T'}^{i})$ deux foncteurs cohomologiques de $C$ dans $C'$,
$\varphi=(\varphi^{i})$ un homomorphisme de foncteurs\pageoriginale
cohomologiques de $T^{\bigdot}$ dans $T'^{\bigdot}$. On suppose 
$T^{\bigdot}$ et $T'^{\bigdot}$ nuls en
degrés $<0$, et $T^{i}$ effaçable pour $i>0$. Soit $n$ un entier. Les
conditions suivantes sont équivalentes:
\begin{enumerate}
\renewcommand{\theenumi}{\alph{enumi}}
\renewcommand{\labelenumi}{\theenumi)}
\item $\varphi^{i}$ est un isomorphisme pour $i\leq n$, un
  monomorphisme pour $i=n+1$.

\item (Si $n\geq -1$). L'homomorphisme $\varphi^{i}$ est un
  monomorphisme si $i=0$, et un épimorphisme si $i\leq n$.

\item (Si $n\geq 0$). L'homomorphisme $\varphi^{\circ}$ est un
  isomorphisme, et les ${T'}^{i}$ sont effaçables pour $0<i\leq n$.
\end{enumerate}

Si ces conditions sont remplies, et si $A\in \Ob(C)$, alors
$\varphi^{n+1}(A):T^{n+1}(A)\rightarrow {T'}^{n+1}(A)$ est un
épimorphisme (donc un isomorphisme) si et seulement si ${T'}^{n+1}(A)$
est effaçable, \ie s'il existe un monomorphisme $A\rightarrow B$ tel que
l'homomorphisme ${T'}^{n+1}(A)\rightarrow {T'}^{n+1}(B)$ correspondant
soit nul. Si $C'=(\Ab)$ (catégorie des groupes abéliens), et si $\xi\in
{T'}^{n+1}(A)$, alors $\xi$ est dans l'image de $T^{n+1}(A)$ si et
seulement si $\xi$ est effaçable, \ie s'il existe un monomorphisme
$A\rightarrow B$ tel que l'image de $\xi$ dans ${T'}^{n+1}(B)$ soit nulle.
\end{enonce*}

\noindent
\emph{Démonstration}. Évidemment a) implique b) et c). Comme les
hypothèses faites sur $T^{\bigdot}$ impliquent que les $T^{i}$ $(i>0)$ sont les
satellites droits de $T^{\circ}$, l'implication c) $\Rightarrow$ b)
résulte de la caractérisation axiomatique des foncteurs satellites;
elle résulte également, par récurrence sur $n$, de la dernière
assertion de 6.2, qui se démontre par l'argument bien connu qu'on se
dispense se répéter ici (\cf démonstration de 6.$''$
ci-dessous). L'implication b) $\Rightarrow$ c) est triviale si $n=1$,
et si $n\geq 0$ on procède par récurrence sur $n$: l'hypothèse de
récurrence nous permet de supposer que $\varphi^{i}$ est un
isomorphisme pour $i\leq n-1$,\pageoriginale un monomorphisme pour
$i=n$, donc par l'hypothèse un isomorphisme aussi pour $i=n$, et il
reste à montrer que c'est un monomorphisme pour $i=n+1$, utilisant que
$\varphi^n$ est un isomorphisme. Soit $A\rightarrow B$ un
monomorphisme qui efface $T^{n+1}$, d'où une suite exacte $0\rightarrow
A\rightarrow B\rightarrow C\rightarrow 0$ dans $C$, et un
homomorphisme de suites exactes dans $C'$:
$$
\xymatrix@R=1.25cm{
T^{n}(B)  \ar[r] \ar[d]& T^{n}(C)  \ar[r]\ar[d] & T^{n+1}(A) \ar[d] \ar[r] & T^{n+1}(B)\ar[d]\\
T'^{n}(B)  \ar[r] & T'^{n}(C) \ar[r] & T'^{n+1}(A)  \ar[r] &
T'^{n+1}(B),
}
$$
où les deux premières flèches verticales sont des isomorphismes
d'après ce qu'on vient de voir, et $T^{n+1}(A)\rightarrow T^{n+1}(B)$
est nul par construction. Alors le lemme des cinq implique que
$T^{n+1}(A)\rightarrow {T'}^{n+1}(A)$ est un monomorphisme, \cqfd

\begin{enonce*}{Corollaire 6.3}\etiquette[6.3]{XII.cor6.3}
Supposons que $T^{\bigdot}$, ${T'^{\bigdot}}$ proviennent de foncteurs
cohomologiques (dénotés par les mêmes symboles) sur la catégorie
dérivée droite $[3]$ $D^{+}(C)$, à valeur dans $C'$; on ne suppose pas
nécessairement les $T^{i}(i>0)$ effaçables. On suppose que
$K^{\bigdot}\in D^{+}(C)$, $H^{i}(K^{\bigdot})=0$ pour $i<0$ implique 
$T^{i}(K^{\bigdot})=T'^{\bigdot}(K^{\bigdot})=0$ pour $i<0$. Soit
$K^{\bigdot}\in \Ob D^{+}(C)$ un complexe borné à gauche dans $C$ tel
que $H^{i}(K^{\bigdot})=0$ pour $i<0$.
\begin{enumerate}
\renewcommand{\theenumi}{\roman{enumi}}
\renewcommand{\labelenumi}{(\theenumi)}
\item Si la condition {\rm a)} de $6.2$ est satisfaite, alors
  l'homomorphisme
  $\varphi^{i}(K^{\bigdot})=T^{i}(K^{\bigdot})\rightarrow
  {T'}^{i}(K^{\bigdot})$ est un isomorphisme pour $i<n$, un
  monomorphisme pour $i=n+1$.

\item Soit de plus $L^{\bigdot}\in \Ob D^{+}(C)$ tel que
  $H^{i}(L^{\bigdot})=0$ pour $i<0$, et soit $u:L^{\bigdot}\rightarrow
  K^{\bigdot}$ un homomorphisme dans $D^{+}(C)$, induisant un
  homomorphisme injectif
  $H^{\circ}(u):H^{\circ}(L^{\bigdot})\rightarrow
  H^{\circ}(K^{\bigdot})$. Alors, pour que
  $T^{n+1}(L^{\bigdot})\rightarrow {T'}^{n+1}(L^{\bigdot})$ soit
  surjectif,\pageoriginale il faut et il suffit que l'image de
  $T^{n+1}(L^{\bigdot})$ dans $T^{n+1}(K^{\bigdot})$ soit contenue dans
  celle de $T^{n+1}(K^{\bigdot})$. Si $C'=(\Ab)$, et si $\xi \in
  {T'}^{n+1}(A)$, alors $\xi$ est dans l'image de $T^{n+1}(A)$ si et
  seulement si son image dans ${T'}^{n+1}(K^{\bigdot})$ est dans celle
  de $T^{n+1}(K^{\bigdot})$.
\end{enumerate}
\end{enonce*}

Rappelons que l'hypothèse que $T^{\bigdot}$ et $T'^{\bigdot}$
proviennent de foncteurs cohomologiques sur $D^{+}(C)$ peut
s'expliciter en disant qu'ils proviennent de foncteurs cohomologiques
sur la catégorie abélienne des complexes bornés à gauche de $C$,
commutant aux translations des degrés, transformant des homomorphismes
de complexes homotopes à zéro en des homomorphismes nuls de $C'$, et
transformant tout homomorphisme $u:K^{\bigdot}\rightarrow L^{\bigdot}$
de complexes qui est un quasi-isomorphisme (\ie qui induit des
isomorphismes $H^{i}(K^{\bigdot})\rightarrow H^{i}(L^{\bigdot})$) en un
isomorphisme de $C'$. Pratiquement, tous les foncteurs cohomologiques
de $C$ dans $C'$ qu'on rencontre sont obtenus ainsi.

\medskip
\noindent
\emph{Démonstration de} \hyperref[XII.cor6.3]{6.3}. Prouvons d'abord (i). L'assertion est
évidente si $H^{i}(K^{\bigdot})=0$ pour $i\leq n+1$, car alors d'après
les propriétés de degrés imposées à $T^{\bigdot}$, $T'^{\bigdot}$, on
aura $T^{i}(K^{\bigdot})={T'}^{i}(K^{\bigdot})=0$ pour $i\leq
n+1$. Soit d'autre part $m$ un entier $\geq 0$ tel que
$H^{i}(K^{\bigdot})=0$ pour $i<m$: il en existe, par exemple $m=0$
fait l'affaire. Nous procédons par récurrence descendante sur $m$,
l'assertion étant prouvée si $m=n+2$. Nous supposons donc l'assertion
prouvée pour les $m'>m$, et la prouvons pour $m$. D'après l'hypothèse
sur les degrés de $K^{\bigdot}$, $K^{\bigdot}$ est quasi-isomorphe à
un complexe à degrés $\geq m$, donc on peut le supposer à degrés $\geq
m$. Il s'ensuit une suite exacte de complexes
$$
0\longrightarrow H^m[-m]\longrightarrow K^{\bigdot}\longrightarrow
K'^{\bigdot}\longrightarrow 0{\quoi},
$$
où $H^m=H^m(K^{\bigdot})$, où le signe $[-m]$ indique translation de
$-m$ sur les degrés d'un complexe,\pageoriginale et où $K'^{\bigdot}$
est défini comme conoyau. Donc les objets de cohomologie de
$K'^{\bigdot}$ sont ceux de $K^{\bigdot}$ en degrés $\geq m+1$, et
nuls en degrés $\leq m+1$; en particulier, on peut appliquer
l'hypothèse de récurrence à $K'^{\bigdot}$. La suite exacte précédente
donne lieu à un homomorphisme de suites exactes
$$
\xymatrix{
T^{i-1} (K'^{\;\bigdot}) \ar[d] \ar[r] & T^{i-m}(H^m)  \ar[r]\ar[d] &
T^{i}(K^{\bigdot}) \ar[r] \ar[d]& T^{i} (K'^{\;\bigdot})  \ar[r] \ar[d]&
T^{i+1-m}(H^m) \ar[d]\\
T'^{i-1}(K'^{\;\bigdot})  \ar[r] & T'^{i-m}(H^m)  \ar[r] &
T'^{i}(K^{\bigdot})  \ar[r] & T'^{i}(K'^{\;\bigdot})  \ar[r] & T'^{i+1-m}(H^m).
}
$$
Faisant alors $i\leq n$ dans le diagramme précédent, et tenant compte
de l'hypothèse de récurrence et de $m\geq 0$, on trouve que les
flèches verticales autour de la flèche verticale médiane sont des
isomorphismes et a fortiori des épimorphismes , et que la flèche
verticale de droite est un monomorphisme, donc par le lemme des cinq
la flèche verticale médiane est un épimorphisme. Faisant $i\leq n+1$,
on trouve de même que les flèches verticales autour de la médiane sont
des monomorphismes, la flèche verticale de gauche étant un
isomorphisme et a fortiori un épimorphisme, d'où on conclut par le
lemme des cinq que la flèche verticale médiane est un monomorphisme,
ce qui établit \hyperref[XII.cor6.3]{6.3} (i).

Pour établir la deuxième assertion \hyperref[XII.cor6.3]{6.3} (ii), on utilise de même la
suite exacte
$$
0\longrightarrow L^{\bigdot}\longrightarrow K^{\bigdot}\longrightarrow
K'^{\bigdot}\longrightarrow 0{\quoi},
$$
où $K'^{\bigdot}$ est défini comme le mapping-cylinder de
$L^{\bigdot}\rightarrow K^{\bigdot}$, donc est encore à cohomologie à
degrés positifs. La « chasse au diagramme »  habituelle dans
$$
\xymatrix{
T^{n}(K^{\bigdot})  \ar[r] \ar[d]& T^{n}(K'^{\;\bigdot})  \ar[r] \ar[d]& T^{n+1}
(L^{\cdot})  \ar[r]\ar[d] & {T^{n+1}(K^{\bigdot})}  \ar[r]
\ar[d]&{T^{n+1}(K'^{\;\bigdot})} \ar[d]\\
T'^{n}(K^{\bigdot}) \ar[r] & T'^{n}(K'^{\;\bigdot})  \ar[r] &
T'^{n+1}(L^{\bigdot}) \ar[r] & T'^{n+1}(K^{\bigdot}) \ar[r] & T'^{n+1}(K'^{\bigdot})
}
$$
donne\pageoriginale alors la dernière assertion de \hyperref[XII.cor6.3]{6.3}, compte tenu du
fait que (en vertu de la première partie déjà démontrée) les deux
premières flèches verticales sont des isomorphismes, les deux
dernières des monomorphismes. Cela achève la démonstration de \hyperref[XII.cor6.3]{6.3}.

\begin{enonce*}{Lemme 6.4}\etiquette[6.4]{XII.lem6.4}
Le théorème $(5.1)$ est vrai si $f$ est fini.
\end{enonce*}

Cela résulte aussitôt en effet de VIII 5.5 et 5.8.

\begin{enonce*}{Proposition 6.5}\etiquette[6.5]{XII.prop6.5}
Soit $h:Y\rightarrow X$ un morphisme de schémas quasi-compacts et
quasi-séparés. Dans {\rm(i)} ci-dessous, $I$ désigne un ensemble donné
ayant au moins deux points, dans {\rm(ii)} et {\rm(iii)}, $\LL$
désigne une partie non-vide de l'ensemble $\PP$ des nombres
premiers, et dans {\rm(iii)} $n$ désigne un entier.
\begin{enumerate}
\renewcommand{\theenumi}{\roman{enumi}}
\renewcommand{\labelenumi}{(\theenumi)}
\item Les conditions suivantes sont équivalentes:
\begin{itemize}
\item[a)] Pour tout faisceau d'ensembles $F$ sur $X$, l'application
canonique
$$
H^{\circ}(X,F)\longrightarrow H^{\circ}(Y,h^{\ast}(F))
$$
est injective (\resp bijective).

\item[b)] Pour tout morphisme fini $X'\rightarrow X$, désignant par
$h^{\circ}:Y'\rightarrow X'$ le morphisme déduit de $h$ par le
changement de base $X'\rightarrow X$, l'application canonique
$$
H^{\circ}(X', I_{X'})\longrightarrow H^{\circ}(Y',I_{Y'})
$$
déduite de $h'$ est injective (\resp bijective).

\item[${\rm b')}$] Avec les notations de {\rm b)}, l'application
$U\mapsto {h'}^{-1}(U)$ de l'ensemble des parties de $Y'$ à la fois
ouvertes et fermées dans l'ensemble des parties de $X'$ à la fois
ouvertes et fermées est injective (\resp bijective).

\item[${\rm b'')}$] (Si\pageoriginale $X$ localement noethérien). 
Avec les notations de {\rm b)}, l'application induite sur 
les ensembles de composantes connexes
$$
\Pi_{\circ}(h'):\Pi_{\circ}(X')\longrightarrow \Pi_{\circ}(Y')
$$
est surjective (\resp bijective), ou encore: si $X'$ est non-vide
(\resp connexe non-vide) il en est de même de $Y'$.
\end{itemize}

De plus, si ces conditions sont satisfaites, alors pour tout faisceau
en groupes $F$ sur $X$, le foncteur $P\mapsto h^{\ast}(P)$ de la
catégorie des torseurs ($=$ fibrés principaux homogènes) sous $F$ dans
la catégorie des torseurs sous $h^{\ast}(F)$ est fidèle (\resp
pleinement fidèle), et a fortiori, dans le cas respé, l'application
canonique
$$
H^1(X,F)\longrightarrow H^1(Y,h^{\ast}(F))
$$
est injective. Enfin, sous les mêmes hypothèses, le foncteur
$X'\mapsto Y'=\fprod{X'}{Y}{X}$ de la catégorie des revêtements étales
de $X$ dans la catégorie des revêtements étales de $Y$ est fidèle
(\resp pleinement fidèle).

\item Les conditions suivantes sont équivalentes:
\begin{itemize}
\item[a)] Pour tout faisceau $F$ de groupes ind-$\LL$-finis sur
$X$, l'application canonique
$$
H^1(X,F)\longrightarrow H^1(Y,h^{\ast}(F))
$$
est bijective pour $i=0$, $1$.

\item[b)] Pour tout morphisme fini $X'\rightarrow X$, désignant par
$h':Y'\rightarrow X'$ le morphisme déduit de $h$ par changement de
base, et pour tout $\LL$-groupe fini ordinaire $G$,
l'application canonique
$$
H^1(X',G_{X'})\longrightarrow H^1(Y',G_{X'})
$$
est bijective pour $i=0$, $1$.

\item[${\rm b')}$] (Si\pageoriginale $\LL=\PP$). Avec
les notations de ${\rm b)}$, le foncteur image inverse par $h'$ induit
une équivalence de la catégorie des revêtements étales de $X'$ avec la
catégorie des revêtements étales de $Y'$.

\item[${\rm b'')}$] (Si $X$ noethérien, $\LL=\PP$). Avec
les notations de {\rm b)}, si $X'$ est non-vide, il en est de même de
$Y'$, et si $y'$ est un point géométrique de $Y'$, $x'$ son image dans
$X'$, l'application canonique
$$
\Pi_i(Y',y')\longrightarrow \Pi_i(X', x')\quav \mbox{ pour }
i=0,1{\quoi},
$$
est bijective.
\end{itemize}

\item Les conditions suivantes sont équivalentes:
\begin{itemize}
\item[a)] Pour tout faisceau $F$ de $\LL$-torsion sur $X$,
l'homomorphisme
$$
H^{i}(X,F)\longrightarrow H^{i}(Y,h^{\ast}(F))
$$
est un isomorphisme si $i\leq n$, un monomorphisme si $i=n+1$.

\item[${\rm a')}$] (Si $n\geq -1$). Pour tout $F$ comme ci-dessus,
l'homomorphisme envisagé est injectif si $i=0$, et surjectif si $i\leq
n$.

\item[b)] (Si $n\geq -1$). Pour tout morphisme fini $X'\rightarrow X$,
désignant par $h':Y'\rightarrow X'$ le morphisme déduit de $h$ par
changement de base $X'\rightarrow X$, et pour tout $\ell\in
\LL$, $\nu>0$, l'homomorphisme canonique
$$
H^{i}(X',(\ZZ/\ell^{\nu}\ZZ_{X'})\longrightarrow
H^{i}(Y',\ZZ/\ell^{\nu}\ZZ)_{Y'})
$$
est injectif pour $i=0$, surjectif pour $i\leq n$.
\end{itemize}
\end{enumerate}
\end{enonce*}

De plus, on a le complément suivant à l'énoncé précédent:

\begin{enonce*}{Corollaire 6.6}\etiquette[6.6]{XII.cor6.6}
\begin{enumerate}
\renewcommand{\theenumi}{\roman{enumi}}
\renewcommand{\labelenumi}{(\theenumi)}
\item Supposons la condition non respée \hyperref[XII.prop6.5]{6.5} {\rm (i) a)} satisfaite,
soit $F$ un faisceau d'ensembles sur $X$, et soit $\xi\in
H^{\circ}(Y,h^{\ast}(F))$. Supposons qu'il\pageoriginale 
existe un monomorphisme de
faisceaux d'ensembles $F\rightarrow G$ tel que l'image de $\xi$ dans
$H^{\circ}(Y,h^{\ast}(G))$ soit dans l'image de
$H^{\circ}(X,G)$. Alors $\xi$ est dans l'image de $H^{\circ}(X,F)$.

\item Supposons la condition {\rm \hyperref[XII.prop6.5]{6.5} (i) a)} respée satisfaite, soit
$F$ un faisceau de groupes sur $X$, et soit $\xi\in
H^1(Y,h^{\ast}(F))$. Supposons qu'il existe un monomorphisme de
faisceaux de groupes $F\rightarrow G$ tel que l'image de $\xi$ dans
$H^1(Y,h^{\ast}(G))$ soit dans l'image de $H^1(X,G)$, alors $\xi$ est
dans l'image de $H^1(X,F)$.

\item Supposons la condition {\rm \hyperref[XII.prop6.5]{6.5} (iii) a)} satisfaite (pour une
valeur donnée $n$), soit $F$ un faisceau de $\LL$-torsion sur
$X$ et soit $\xi\in H^{n+1}(Y,h^{\ast}(F))$. Supposons qu'il existe
un monomorphisme $F\rightarrow G$ de faisceaux de
$\LL$-torsion, tel que l'image de $\xi$ dans $H^{n+1}
(Y,h^{\ast}(G))$ soit dans l'image de $H^{n+1}(X,G)$, alors $\xi$
est dans l'image de $H^{n+1}(X,F)$.
\end{enumerate}
\end{enonce*}

\noindent
\emph{Démonstration de \hyperref[XII.prop6.5]{6.5} et \hyperref[XII.cor6.6]{6.6}.} 
Prouvons d'abord \hyperref[XII.cor6.6]{6.6}. Le cas
de \hyperref[XII.cor6.6]{6.6} (iii) est un cas particulier de la dernière assertion dans \hyperref[XII.cor6.3]{6.3},
appliqué au cas où $C$ est la catégorie des faisceaux abéliens de
$\LL$-torsion sur $X$, $C'$ est la catégorie $(\Ab)$, et
$T^{i}(F)=H^{i}(X,F)$, ${T'}^{i}(F)=H^{i}(Y,g^{\ast}(F))$ (noter que
$T'^{\bigdot}$ est bien un foncteur cohomologique en $F$, grâce au
fait que le foncteur $h^{\bigdot}$ est exact). Bien entendu, ici la
considération des complexes comme dans \hyperref[XII.cor6.3]{6.3} est inutile (elle nous sera
commode plus loin (\hyperref[XII.prop6.8]{6.8})) et nous utilisons uniquement la dernière
assertion de \hyperref[XII.cor6.3]{6.3}.

Pour prouver (i), considérons la somme amalgamée $H=G\amalg_F G$,
limite inductive, dans la catégorie des faisceaux d'ensembles sur $X$,
du diagramme
$$
\xymatrix@R=.5cm{
  & G\\
F \ar[ru] \ar[rd]&\\
  & G
}
$$
Les\pageoriginale propriétés d'exactitude habituelles du topos des
faisceaux d'ensembles montrent, comme $F\rightarrow G$ est un
monomorphisme, que l'on a un diagramme \emph{exact} de faisceaux
d'ensembles sur $X$
$$
F\to G\rightrightarrows H
$$
(ce que l'on peut exprimer en disant que dans un topos, tout
monomorphisme est effectif, et prouver en se ramenant comme d'habitude
à la catégorie des ensembles). Comme les foncteurs $H^{\circ}$ et
$h^{\ast}$ sont exacts à gauche, on en conclut un homomorphisme de
diagrammes exacts d'ensembles
$$
\xymatrix@R=1.25cm{
H^{0}(X,F) \ar[r]\ar[d]& H^{0}(X,G)  \ar[d]\ar@<2pt>[r]\ar@<-2pt>[r]& H^{0}(X,H)\ar[d]\\
H^{0}(Y,h^{\ast}(F)) \ar[r]& H^{0}(Y,h^{\ast}(G)) \ar@<2pt>[r]\ar@<-2pt>[r]& H^{0}(Y,h^{\ast}(H)),
}
$$
où les flèches verticales sont injectives. Une « chasse de
diagramme » immédiate prouve alors l'assertion \hyperref[XII.cor6.6]{6.6} (i).

Pour prouver \hyperref[XII.cor6.6]{6.6} (ii), nous allons d'abord prouver la deuxième
assertion de \hyperref[XII.prop6.5]{6.5} (i), savoir que moyennant la condition a) de \hyperref[XII.prop6.5]{6.5} (i),
pour tout faisceau de groupes $F$ sur $X$, le foncteur $P\mapsto
h^{\ast}(P)$, de la catégorie des torseurs sous $F$ dans la catégorie
des torseurs sous $h^{\ast}(F)$, est fidèle (\resp pleinement fidèle),
donc, dans le cas respé, induit une injection sur les $H^1$. Pour
ceci, si $P$ et $P'$ sont deux torseurs sous $F$, désignons par
$\sheaf{\Isom}_F(P,P')$ le faisceau des $F$-isomorphismes de $P$ sur
$P'$. Il est immédiat que la formation de ce faisceau commute à toute
extension de la base, en particulier à l'extension de la base par
$h$. (C'est là une assertion valable pour tout morphisme de
sites). Utilisant l'hypothèse sur l'effet de $h^{\ast}$ sur les
$H^{\circ}$ pour le faisceau $\sheaf{\Isom}_F(P,P')$, on trouve que
l'application $\Isom_F(P,P')\rightarrow
\Isom_F(h^{\ast}(P),h^{\ast}(P'))$ est injective (\resp bijective), ce
qui signifie que le foncteur $P\mapsto h^{\ast}(P)$ est\pageoriginale 
fidèle (\resp
pleinement fidèle). (NB dans la catégorie des torseurs sous un
faisceau en groupes, tout homomorphisme est un isomorphisme). Notons
d'autre part, que si $F$ est un sous-faisceau en groupes d'un faisceau
en groupes $G$, alors la donnée d'un torseur $P$ sous $F$ revient à la
donnée d'un torseur sous $G$ (savoir celui déduit de $P$ par extension
$F\rightarrow G$ du Groupe structural), \emph{muni} d'une section de
$Q/F$ (qui s'interprète en effet comme une « restriction du Groupe
structural »  de $G$ à $F$). Ceci posé, avec les notations de 6.6
(ii), comme l'image de $\eta$ de $\xi$ dans $H^1(Y,h^{\ast}(G))$ provient d'un
élément de $H^1(X,G)$, ce dernier est défini par un torseur $Q$ sous
$G$, de sorte que $\eta$ est la classe du torseur $h^{\ast}(Q)$ sous
$h^{\ast}(G)$. Le fait que $\eta$ provient d'un $\xi\in
H^1(Y,h^{\ast}(F))$ s'explicite alors en disant que $\xi$ est la
classe du $h^{\ast}(F)$-torseur défini par $h^{\ast}(Q)$ et une
section convenable de $h^{\ast}(Q)/h^{\ast}(F)$. Comme le foncteur
$h^{\ast}$ est exact à droite, ce dernier faisceau n'est autre que
$h^{\ast}(Q/F)$, et comme $h^{\ast}$ induit une bijection sur les
$H^{\circ}$ des faisceaux d'ensembles, il s'ensuit que la section
envisagée de $h^{\ast}(Q)/h^{\ast}(F)$ provient d'une section de
$Q/F$. Cette dernière définit alors un $F$-torseur par restriction du
groupe structural, et la classe de ce torseur dans $H^1(X,F)$ a
évidemment $\xi$ comme image dans $H^1(Y,h\xi (F))$. Cela prouve \hyperref[XII.cor6.6]{6.6}
(ii) et achève la démonstration de \hyperref[XII.cor6.6]{6.6}.

Prouvons maintenant \hyperref[XII.prop6.5]{6.5}, en commençant par \hyperref[XII.prop6.5]{6.5} (iii). Évidemment a)
implique ${\rm a'}$), et l'implication inverse est un cas particulier
de \hyperref[XII.lem6.2]{6.2} (appliqué encore au cas où $C$ est la catégorie des faisceaux
de $\LL$-torsion sur $X$, et $C'$ la catégorie des faisceaux de
groupes abéliens), pourvu qu'on prouve l'effaçabilité dans $C$ des
foncteurs $H^{i}(X,F)$ pour $i>0$. Cela n'offre en effet pas de
difficulté, mais nous pouvons aussi nous dispenser de prouver ce
résultat, en\pageoriginale notant que $F$ est limite inductive
filtrante de ses sous-faisceaux $F_i$ pour lesquels il existe un
entier $n>0$ annulant $F_i$ (\hyperrefext[IX][IX.def1.1]{1.1}). Alors $h^{\ast}(F)$ est limite
inductive des $h^{\ast}(F_i)$, et compte tenu de la commutation de la
cohomologie à la limite inductive (VII 3.3), résultant des hypothèses
de quasi-compacité et de quasi-séparation faites sur $X$ et $Y$ (qui
n'avaient pas servi pour \hyperref[XII.cor6.6]{6.6}), on se ramène à prouver la condition a)
pour les $F_i$, ce qui nous ramène au cas où il existe un entier $m>0$
tel que $mF=0$. On peut évidemment supposer que les diviseurs premiers
de $m$ sont $\in \LL$, et on peut alors prendre pour $C$ la
catégorie des faisceaux de $\ZZ/m\ZZ$-Modules sur
$X$. Or on sait que dans cette catégorie, les foncteurs $H^{i}(X, -)$
pour $i\geq 1$ sont effaçables (\quad). Cela montre l'équivalence
des conditions a) et ${\rm a'}$) de \hyperref[XII.prop6.5]{6.5} (iii). D'autre part,
appliquant a) à l'image directe d'un faisceau de $\LL$-torsion
$F'$ sur $X'$, $X'$ fini sur $X$, qui est de $\LL$-torsion
grâce à \hyperrefext[IX][IX.prop1.2]{1.2} (v), \hyperrefext[IX][IX.prop1.6]{1.6} (iii), et 
utilisant \hyperref[XII.lem6.4]{6.4} et VIII 5.5, on trouve
que l'homomorphisme
$$
H^{i}(X',F')\longrightarrow H^{i}(Y',{h'}^{\ast}(F'))
$$
est un isomorphisme pour $i\leq n$, un monomorphisme pour $i=n+1$,
d'où a fortiori la condition b). Reste à prouver que cette condition
implique ${\rm a'}$). Pour ceci, utilisons le fait (\hyperrefext[IX][IX.cor2.7.2]{2.7.2}) que tout
faisceau de $\LL$-torsion $F$ sur $X$ est limite inductive
filtrante de faisceaux de $\LL$-torsion constructibles; alors
l'argument de passage à la limite déjà utilisé montre que dans ${\rm
a'}$), on peut se limiter au cas où $F$ est constructible. En vertu de
\hyperrefext[IX][IX.prop2.14]{2.14} il existe alors un nombre fini de morphismes finis
$p_i:X'_i\rightarrow X$, de faisceaux constants de
$\LL$-torsion $G_{X'_i}$ sur les $X'_i$, et un homomorphisme
injectif
$$
F\longrightarrow \prod_{i}p_{i_{\ast}}(G_{i X'_i})
$$
Ceci\pageoriginale dit, si $n=-1$, \ie lorsque la conclusion voulue se
réduit à l'assertion d'injectivité pour $H^{\circ}(X,F)\rightarrow
H^{\circ}(Y,h^{\ast}(F))$, on voit aussitôt qu'on peut dans cette
question remplacer $F$ par un faisceau $G$ tel qu'il existe un
monomorphisme $F\rightarrow G$, et par suite on peut remplacer $F$ par
le produit des $p_{i^{\ast}}(G_{i X'_i})$. Compte tenu de \hyperref[XII.lem6.4]{6.4}, on est
donc ramené à prouver l'assertion d'injectivité en remplaçant
$h:Y\rightarrow X$ par $h':Y'\rightarrow X'$, et $F$ par un faisceau
constant sur $X'$ de valeur un groupe fini de $\LL$-torsion
ordinaire $G$. Or $G$ est isomorphe à une somme de groupes de la forme
$\ZZ/\ell^{\nu}\ZZ$, avec $\ell\in \LL$. On est
donc bien réduit à vérifier b). Lorsque $n\geq 0$, on procède de fa\c
con analogue. Par récurrence sur $n$, on peut supposer la conclusion
voulue (et par suite a)) prouvée pour l'entier $n-1$ au lieu de
$n$. Utilisant alors \hyperref[XII.cor6.6]{6.6} (iii), on voit qu'on peut encore remplacer
$F$ par un faisceau $G$ tel qu'il existe un monomorphisme
$F\rightarrow G$, et on prendra encore pour $G$ le faisceau
$p_{\ast}(G_{X'})$. Compte tenu de \hyperref[XII.lem6.4]{6.4} et de la nullité des
$R^{i}p_{\ast}$ $(i\geq 1)$ pour un morphisme fini (VIII 5.5), cela
nous ramène encore à prouver la surjectivité pour $H^n$ en rempla\c
cant $h:Y\rightarrow X$ par un $h':Y'\rightarrow X'$, et $F$ par un
faisceau constant sur $X'$ de valeur un groupe fini de
$\LL$-torsion ordinaire. Comme l'injectivité est déjà acquise
grâce à l'hypothèse de récurrence et l'implication ${\rm
a')\Rightarrow a)}$, on est encore ramené, par dévissage sur $G$, au
cas où $G$ est de la forme $\ZZ/\ell^{\nu}\ZZ$,
c'est-à-dire au cas envisagé dans la condition b) de \hyperref[XII.prop6.5]{6.5} (iii). Cela
prouve le cas (iii) de \hyperref[XII.prop6.5]{6.5}.

Les cas (i) et (ii) se démontrent exactement de la même façon, pour
ce qui est de l'équivalence des conditions a) et b) dans
\hyperref[XII.prop6.5]{6.5} (i) \resp 
\hyperref[XII.prop6.5]{6.5} (ii): on\pageoriginale utilise \hyperrefext[IX][IX.cor2.7.2]{2.7.2} et
la commutation du 
$H^{\circ}$ et $H^1$ aux limites inductives filtrantes de faisceaux
(VII 5.14) pour se ramener au cas où $F$ est constructible, puis 
\hyperrefext[IX][IX.prop2.14]{2.14}, \hyperref[XII.cor6.6]{6.6} (i) (ii) et VIII 5.8 pour 
se ramener au cas particulier
envisagé dans b). Par ailleurs, dans (i), l'équivalence des conditions
b), ${\rm b'}$) et ${\rm b''}$) est triviale et n'est mise que pour
mémoire, ainsi que le fait qu'elles impliquent la fidélité (\resp la
pleine fidélité) du foncteur $X'\mapsto Y'$ de la catégorie des
revêtements étales de $X$ dans la catégorie des revêtements étales de
$Y$. Il en est de même de l'équivalence des conditions b), ${\rm b'}$)
et ${\rm b''}$) dans \hyperref[XII.prop6.5]{6.5} (ii), qui ont été ajoutées pour
faire bon 
poids. La démonstration de \hyperref[XII.prop6.5]{6.5} est ainsi achevée.

\begin{enonce*}[remark]{Remarque 6.7}\etiquette[6.7]{XII.rem6.7}
\footnote{Le rédacteur recommande d'omettre la lecture de ces
remarques, ainsi que de (\hyperref[XII.rems6.13]{6.13}), introduites subrepticement (en même
temps que diverses autres modifications plus ou moins heureuses du
texte original) par un collaborateur irrévérencieux.}.
\begin{enumerate}
\renewcommand{\labelenumi}{\theenumi)}
\item Évidemment, les arguments démontrant \hyperref[XII.prop6.5]{6.5} et
  \hyperref[XII.cor6.6]{6.6} sont de nature 
très générale et essentiellement triviale et auraient intérêt à être
dégagés en des lemmes abstraits de la même eau que \hyperref[XII.lem6.2]{6.2}
(les faisceaux 
$p_{\ast}(G_{X'})$ jouant le rôle de cogénérateurs de la catégorie $C$
dans laquelle on travaille). On laisse ce plaisant exercice au
lecteur, et nous nous bornons à signaler que le même énoncé
essentiellement pourrait être donné en regardant des foncteurs tels
que $R^{i}f_{\ast}$ au lieu des foncteurs $H^{i}$. Les mêmes remarques
s'appliquent à \hyperref[XII.prop6.8]{6.8}, \hyperref[XII.prop6.11]{6.11} ci-dessous.

\item Lorsque $X$ est noethérien, alors les arguments donnés montrent
que dans l'énoncé des conditions b) et leurs variantes dans \hyperref[XII.prop6.5]{6.5} (i)
(ii) (iii), on peut se borner à prendre $X'$ \emph{intègre}; lorsque
$X$ est universellement japonais, on peut même les prendre intègre et
normaux; c'est également possible sans la\pageoriginale restriction
japonaise sur $X$, à condition de prendre des morphismes $p:X'\rightarrow X$
\emph{entiers} au lieu de morphismes finis. Nous ne nous servirons
d'aucune de ces variantes dans la suite du séminaire.

\item La démonstration de \hyperref[XII.cor6.6]{6.6} n'utilise pas l'hypothèse de
quasi-compacité et de quasi-séparation faite sur $X$, et il en est de
même pour le fait que dans (i) (ii) (iii) la condition a) implique les
autres, ainsi que pour l'équivalence des conditions b), ${\rm b'}$),
${\rm b''}$) entre elles. Notons d'autre part que la forme b) des
conditions 6.5 (i) montre que, pour $\LL=\phi$, ce sont des cas
particuliers (lorsque $X$ est quasi-compact quasi-séparé) de \hyperref[XII.prop6.5]{6.5}
(iii), pour $n=-1$ \resp pour $n=0$. Cela reste d'ailleurs vrai (en
se bornant aux conditions sous la forme a)) sans la restriction
énoncée sur $X$, comme on voit aisément grâce au fait que pour tout
faisceau d'ensembles $F$, et tout nombre premier $\ell$, on peut
trouver deux faisceaux abéliens de $\ell$-torsion $G$ et $H$, et un
homomorphisme $u:G\rightarrow H$ des faisceaux d'ensembles
sous-jacents, tels que $F$ soit isomorphe à l'image inverse de la
section nulle (on prendra pour $F$ le $\ZZ/\ell
\ZZ$-Module libre engendré par $F$, et pour $H$ le faisceau
constant $\ZZ/\ell\ZZ$ sur $X$). Le même argument montre
que (sans restriction de quasi-compacité et de quasi-séparation) \hyperref[XII.prop6.5]{6.5}
(ii) a) (où il suffit même de faire $i=0$) implique la condition
respée \hyperref[XII.prop6.5]{6.5} (i) a); cela permet par suite, compte tenu que cette
dernière implique déjà l'injectivité de $H^1(X,F)\rightarrow
H^1(Y,h^{\ast}(F))$ pour tout faisceau en groupes $F$ sur $X$, de se
borner dans l'énoncé de \hyperref[XII.prop6.5]{6.5} (ii) a) d'exiger pour $i=1$ la
\emph{surjectivité} de $H^1(X,F)\rightarrow
H^1(Y,h^{\ast}(F))$. Malheureusement, le cas \hyperref[XII.prop6.5]{6.5} (ii) ne peut être
envisagé comme cas particulier de \hyperref[XII.prop6.5]{6.5} (iii), ce qui nous oblige
souvent de répéter dans le cas non commutatif un argument déjà fait
pour l'essentiel dans le cas commutatif. Notons également que, sauf
\hyperref[XII.cor6.10]{6.10}, les résultats qui suivent sont énoncés de sorte qu'ils sont valables
sans\pageoriginale hypothèse de quasi-compacité et de quasi-séparation.
\end{enumerate}
\end{enonce*}

\begin{enonce*}{Proposition 6.8}\etiquette[6.8]{XII.prop6.8}
(Lemme de descente). Avec les notations générales de \hyperref[XII.prop6.5]{6.5}, supposons
donné de plus un morphisme $p:\bar{X}\rightarrow X$, désignons par
$\bar{h}:\bar{Y}\rightarrow \bar{X}$ le morphisme déduit de $h$ par le
changement de base $p$, et par $p_y:\bar{Y}\rightarrow Y$ le morphisme
canonique. Nous supposons, ou bien que $p$ est surjectif, ou bien que
$h$ est une immersion fermée et que $X=p(\bar{X})\cup h(Y)$.
\begin{enumerate}
\renewcommand{\theenumi}{\roman{enumi}}
\renewcommand{\labelenumi}{(\theenumi)}
\item Supposons que le morphisme $\bar{h}:\bar{Y}\rightarrow \bar{X}$
satisfasse à la condition d'injectivité (\resp de bijectivité) de
\hyperref[XII.prop6.5]{6.5} {\rm(i) a)}. Dans le cas respé, supposons de plus que le
morphisme de changement de base
\begin{equation*}
h^{\ast}(p_{\ast}(\bar{F}))\longrightarrow
p_{Y_{\ast}}(\bar{h}^{\ast}(\bar{F}))\tag{$\ast$} 
\end{equation*}
est injectif pour tout faisceau d'ensembles $\bar{F}$ sur
$\bar{X}$. Alors le morphisme $h:X\rightarrow Y$ satisfait également
la condition d'injectivité (\resp de bijectivité) de $6.5$ {\rm(i)
a).}

\item Supposons que $\bar{h}:\bar{Y}\rightarrow \bar{X}$ satisfasse la
condition de $6.5$ {\rm(ii) a)}, et que pour tout faisceau de
{\rm ind}-$\LL$-groupes $\bar{F}$ sur $\bar{X}$, le morphisme de
changement de base $(\ast)$ soit bijectif. Alors $h:Y\rightarrow X$
satisfait également la condition de {\rm \hyperref[XII.prop6.5]{6.5} (ii) a).}

\item Supposons que $\bar{h}:\bar{Y}\rightarrow \bar{X}$ satisfasse à
la condition de {\rm \hyperref[XII.prop6.5]{6.5} (iii) a),} et que pour tout faisceau abélien
$\bar{F}$ de $\LL$-torsion sur $\bar{X}$, l'homomorphisme de
changement de base
$$
h^{\ast}(R^{i}p_{\ast}(\bar{F}))\longrightarrow
R^{i}p_{Y_{\ast}}(\bar{h}^{\ast}(\bar{F}))
$$
soit bijectif pour $i\leq n-1$, injectif pour $i=n$ (où $n$ est un
entier $\geq 1$ donné). Alors le morphisme $h:Y\rightarrow X$
satisfait également à la condition de {\rm \hyperref[XII.prop6.5]{6.5} (iii) a).}
\end{enumerate}
\end{enonce*}

\noindent
\emph{Démonstration de} \hyperref[XII.prop6.8]{6.8}.\pageoriginale Dans le cas où on ne
suppose pas $p$ surjectif, mais $h$ une immersion fermée et
$X=p(\bar{X})\cup h(Y)$, nous considérons le schéma somme $\bar{X}'$
de $\bar{X}$ et $Y$, et le morphisme $p':\bar{X}'\rightarrow X$ déduit
de $p$ et $h$. Alors $p'$ est surjectif, de plus on vérifie
trivialement, dans chacun des trois cas envisagés (i) (ii) (iii), que
les hypothèses faites sur le couple $(h,p)$ sont encore satisfaites
pour le couple $(h,p')$. Cela nous ramène donc au cas où $p$ est
surjectif, dans lequel nous allons nous placer dans la suite. Dans ce
cas pour tout faisceau $F$ sur $X$, l'homomorphisme canonique
$$
F\longrightarrow p_{\ast}(p^{\ast}(F))
$$
est injectif. Dans le cas (i), pour vérifier l'injectivité de
$H^{\circ}(X,F)\rightarrow H^{\circ}(Y,h^{\ast}(F))$, on est ramené
aussitôt au cas où on remplace $F$ par $p_{\ast}(\bar{F})$ (où on pose
$\bar{F}=p^{\ast}(F)$). Or considérons le diagramme commutatif
$$
\xymatrix@C=2.75cm@R=.75cm{
**[r]H^{\circ}(X,p_{\ast}(\bar{F})) \ar[dd]^{\text{\rotatebox{90}{$\backsim$}}} \ar[r]& **[r]H^{\circ}(Y,h^{\ast}(\bar{F})) \ar[d]\\
                           &**[r]H^{\circ}(Y,p_{Y^{\ast}}(\bar{h}^{\ast}(\bar{F})))\ar[d]_{\text{\rotatebox{90}{$\backsim$}}}\\
**[r]H^{\circ}(\bar{X},\bar{F})  \ar[r] &
                           **[r]H^{\circ}(\bar{Y},\bar{h}^{\ast}(\bar{F})),  
}
$$
où la flèche verticale gauche et la deuxième flèche verticale droite
sont les isomorphismes canoniques, la première flèche verticale droite
provenant de l'homomorphisme de changement de base par $h$. Par
hypothèse la deuxième flèche horizontale est injective, ce qui
implique aussitôt qu'il en est de même de la première, ce qui prouve
l'assertion non respée de (i). Pour l'assertion respée, on suppose que
la deuxième flèche horizontale est surjective, de plus l'homomorphisme
de changement de base étant un monomorphisme par hypothèse, il en est
de même de la première flèche verticale droite, donc du
composé\pageoriginale des deux flèches verticales droites. Il s'ensuit
aussitôt que la première flèche horizontale est également surjective,
ce qui achève de prouver (i). Pour prouver (ii), grâce à ce qui
précède et \hyperref[XII.prop6.5]{6.5} (i) on est ramené à prouver la surjectivité pour
$H^1(X,F)\rightarrow H^1(Y,h^{\ast}(F))$. Utilisant \hyperref[XII.cor6.6]{6.6} (ii) on est
encore ramené à la prouver pour $F$ de la forme $p_{\ast}(\bar{F})$,
où $\bar{F}$ est un faisceau de ind-$\LL$-groupes sur $\bar{X}$
(\hyperrefext[IX][IX.prop1.6]{1.6} (ii)). Écrivant les suites exactes pour les $H^1$, (variantes
de (\hyperref[XII.prop3.2]{3.2})) pour les morphismes $p$ et $p_Y$, on trouve un homomorphisme
de suites exactes non commutatives:
$$
\xymatrix@R=1.5cm{
H^{1} (X,p_{\ast}(\bar{F}))  \ar[r] \ar[d] & H^{1}(\bar{X},\bar{F})  \ar[r]\ar[d] &
H^{0}(X,R^{1}p_{\ast}(\bar{F})) \ar[d]&\\
H^{1}(Y,p_{Y^{\ast}}(\bar{h}^{\ast}(\bar{F})) \ar[r] &
H^{1}(\bar{Y},\bar{h}^{\ast}(\bar{F}))  \ar[r] &
H^{0}(Y,R^{1}(p_{Y^{\ast}}(\bar{h}^{\ast}(\bar{F})))), 
}
$$
où dans chaque ligne la première flèche est injective et identifie le
premier terme à l'image inverse, par la deuxième flèche, du point
marqué du troisième terme. Par hypothèse sur $\bar{h}$, la flèche
verticale médiane est bijective. D'autre part, les deux flèches
verticales extrêmes se factorisent respectivement par
$H^1(Y,h^{\ast}(p_{\ast}(\bar{F}))$ et
$H^{\circ}(Y,h^{\ast}(R^1(p_{\ast}(\bar{F}))$, en composant les
homomorphismes $H^{i}(X,-)\rightarrow H^1(Y,-)$ (pour $i=1$ et $i=0$
respectivement) et les homomorphismes déduits en appliquant
$H^{i}(Y,-)$ (pour ces mêmes valeurs de i) à l'homomorphisme de
changement de base $h^{\ast}(p_{\ast}(\bar{F}))\rightarrow
p_{Y^{\ast}}(\bar{h}^{\ast}(\bar{F}))$ \resp
$h^{\ast}(R^1p_{\ast}(\bar{F}))\rightarrow
R^1p_{Y^{\ast}}(\bar{h}^{\ast}(\bar{F}))$. Comme la première de ces
flèches est un isomorphisme par hypothèse, il s'ensuit que la première
flèche verticale du diagramme ci-dessus s'identifie à l'homomorphisme
$H^1(C,-)\rightarrow H^1(Y,-)$ que nous voulons étudier. Pour prouver
sa bijectivité, il reste donc à prouver que la dernière flèche
verticale du diagramme est injective. Or d'après ce qu'on vient de
dire, elle est composée de deux applications,\pageoriginale dont la
première est un homomorphisme canonique $H^{\circ}(X,-)\rightarrow
H^{\circ}(Y,-)$, donc bijective, et la deuxième est injective, car
déduite par le foncteur exact à gauche $H^{\circ}(Y,-)$ d'un
homomorphisme $H^{\ast}(R^1(p_{\ast}(\bar{F}))\rightarrow
R^1p_{\ast}(\bar{h}^{\ast}(\bar{F}))$ qui est un monomorphisme. Ce
dernier fait résulte en effet aisément de l'hypothèse de bijectivité
faite pour l'homomorphisme de changement de base dans la dimension
précédente 0, par exemple en utilisant la deuxième assertion de \hyperref[XII.prop6.5]{6.5}
(i) et le calcul habituel des fibres des images directes supérieures
(VIII 5.3). Cela prouve \hyperref[XII.prop6.8]{6.8} (ii).

Reste à prouver \hyperref[XII.prop6.8]{6.8} (iii). Pour ceci, il nous sera plus commode, au
lieu d'utiliser directement l'injection $F\rightarrow
p_{\ast}(p^{\ast}(F))$ comme ci-dessus, d'utiliser l'homomorphisme
$$
F\longrightarrow Rp_{\ast}(\bar{F})
$$
dans la catégorie dérivée droite $D^{+}(C)$, où $C$ désigne la
catégorie des faisceaux abéliens sur $X$. Rappelons que par
définition, $Rp_{\ast}(\bar{F})$ est le complexe
$p_{\ast}(C(\bar{F}))$, où $C(\bar{F})$ est un complexe résolution
injective de $\bar{F}$ dans la catégorie des faisceaux abéliens sur
$\bar{X}$. Si $K^{\bigdot}$ est un complexe sur $X$, nous désignons
par $H^{i}(X,K^{\bigdot})$ les groupes d'hypercohomologie de $X$ à
coefficients dans $K^{\bigdot}$, et on utilise les notations analogues
sur $Y$. En vertu de \hyperref[XII.cor6.3]{6.3}, on est ramené à prouver la surjectivité pour
\begin{equation*}
H^n(X,Rp_{\ast}(\bar{F}))\longrightarrow
H^n(Y,h^{\ast}Rp_{\ast}(\bar{F})){\quoi}.\tag{1}
\end{equation*}\etiquette[1]{XII.eq1}
Or par hypothèse, l'homomorphisme de changement de base
\begin{equation*}
H^{\ast}(Rp_{\ast}(\bar{F}))\longrightarrow Rp_{Y\ast}(\bar{h}^{\ast}(\bar{F}))\tag{2}
\end{equation*}\etiquette[2]{XII.eq2}
induit un isomorphisme sur les faisceaux de cohomologie en degrés
$\leq n-1$, et un monomorphisme en degrés $n$, ce qui peut s'exprimer
en disant que le « mapping-cylinder » \pageoriginale $K^{\bigdot}$
de l'homomorphisme précédent n'a des faisceaux de cohomologie non nuls
qu'en degrés $\geq n$. Donc $H^{i}(Y,K^{\bigdot})=0$ si $i<n$, ce qui
implique, par la suite exacte de cohomologie, que les homomorphismes
\begin{equation*}
H^{i}(Y,h^{\ast}(Rp_{\ast}(\bar{F}))\longrightarrow H^{i}(Y,Rp_{Y\ast}(\bar{h}^{\ast}(\bar{F})){\quoi},\tag{3}
\end{equation*}\etiquette[3]{XII.eq3}
induits par l'homomorphisme ci-dessus, sont des isomorphismes pour
$i\leq n-1$, un monomorphisme pour $i=n$. Considérons le composé de
(1) et (3)
\begin{equation*}
H^n(X,Rp_{\ast}(\bar{F}))\longrightarrow
H^n(Y,h^{\ast}(Rp_{\ast}(\bar{F})))\longrightarrow H^n(Y,Rp_{Y\ast}(\bar{h}^{\ast}(\bar{F})){\quoi}.\tag{4}
\end{equation*}\etiquette[4]{XII.eq4}
Les termes extrêmes sont respectivement isomorphes, en vertu des
définitions, à $H^n(\bar{X},\bar{F})$ et
$H^n(\bar{Y},\bar{h}^{\ast}(\bar{F}))$, et le composé de (4) n'est
autre que l'homomorphisme déduit de $\bar{h}$, qui est bijectif grâce
à l'hypothèse faite sur $\bar{h}$. Comme la deuxième flèche de (4) est
injective d'après ce qu'on vient de voir, la première est également
bijective, ce qui achève la démonstration de \hyperref[XII.prop6.8]{6.8}.

D'ailleurs, on constate aussitôt que la démonstration précédente
fournit le résultat suivant, légèrement plus précis et plus général:

\begin{enonce*}{Corollaire 6.9}\etiquette[6.9]{XII.cor6.9}
Les notations sont celles de \hyperref[XII.prop6.8]{6.8}. On suppose de plus satisfaite pour
$h$ la condition non respée de \hyperref[XII.prop6.5]{6.5} {\rm(i)} (\resp la condition
respée de \hyperref[XII.prop6.5]{6.5} {\rm(i)}, \resp la condition de {\rm \hyperref[XII.prop6.5]{6.5} 
(iii)}, avec
$n\geq -1$), et que pour tout faisceau d'ensembles $\bar{F}$ sur
$\bar{X}$, l'homomorphisme de changement de base
$h^{\ast}(p_{\ast}(\bar{F}))\rightarrow
p_{Y^{\ast}}(\bar{h}^{\ast}(\bar{F}))$ est injectif (\resp bijectif,
\resp que pour tout faisceau de $\LL$-torsion $\bar{F}$ sur
$\bar{X}$, l'homomorphisme de changement de base
$h^{\ast}R^{i}p_{\ast}(\bar{F})\rightarrow
R^{i}p_{Y^{\ast}}(\bar{h}^{\ast}(\bar{F}))$ soit bijectif pour $i\leq
n$, injectif pour $i=n+1$). Soit $F$ un faisceau d'ensembles (\resp de
ind-$\LL$-groupes, \resp un faisceau de $\LL$-torsion)
sur $X$, et soit $\xi \in H^{n+1}(Y,h^{\ast}(F))$ (où on prend $n=-1$
dans le cas\pageoriginale {\rm(i)}, $n=0$ dans le cas {\rm(ii)}). Posons
$\bar{F}=p^{\ast}(F)$. Alors, pour que $\xi$ soit dans l'image de
$H^{n+1}(X,F)$, il faut et il suffit que son image inverse dans
$H^{n+1}(\bar{Y},p_{Y^{\ast}}(h^{\ast}(F))=H^{n+1}(\bar{Y},\bar{h}^{\ast}(\bar{F}))$ 
soit dans l'image de $H^{n+1}(\bar{X},\bar{F})$.
\end{enonce*}

Nous utiliserons cette forme précisée de \hyperref[XII.prop6.8]{6.8} pour prouver le

\begin{enonce*}{Corollaire 6.10}\etiquette[6.10]{XII.cor6.10}
Les notations sont celles de \hyperref[XII.prop6.5]{6.5}. On suppose, pour simplifier,
vérifiée la condition non respée de {\rm \hyperref[XII.prop6.5]{6.5} (i)}.
\begin{enumerate}
\renewcommand{\theenumi}{\roman{enumi}}
\renewcommand{\labelenumi}{(\theenumi)}
\item Pour que la condition respée de {\rm \hyperref[XII.prop6.5]{6.5} (i)} soit satisfaite
(\resp pour qu'on ait {\rm \hyperref[XII.prop6.5]{6.5} (ii) a),} \resp pour qu'on ait {\rm \hyperref[XII.prop6.5]{6.5}
(iii) a))}, il faut et il suffit que la condition suivante soit
satisfaite: Pour tout faisceau $F$ d'ensembles (\resp de
ind-$\LL$-groupes, \resp de groupes abéliens de
$\LL$-torsion) sur $X$, tout $\xi\in H^{\circ}(Y,h^{\ast}(F))$
(\resp tout $\xi\in H^{i}(Y,h^{\ast}(F))$ avec $i=0,1$, \resp tout
$\xi \in H^{i}(Y,h^{\ast}(F))$ avec $i\leq n$), et toute partie fermée
non vide $X'$ de $X$, désignant par $Y'$ son image inverse dans $Y$,
il existe un morphisme $p:\bar{X}\rightarrow X$ satisfaisant les
conditions suivantes:
\begin{itemize}
\item[$1^{\circ}$)] L'image $p(\bar{X})$ est contenue dans $X'$ et
contient un ouvert non vide de $X'$.

\item[$2^{\circ}$)] Pour tout faisceau d'ensembles $\bar{F}$ sur
$\bar{X}$ (\resp\ldots) l'homomorphisme de changement de base
$h^{\ast}(p_{\ast}(\bar{F}))\rightarrow p_{Y^{\ast}}(\bar{h}^{\ast}(\bar{F}))$ 
est un isomorphisme (\resp pour
tout faisceau de ind-$\LL$-groupes $\bar{F}$ sur $X$,
l'homomorphisme de changement de base
$h^{\ast}(R^{i}p_{\ast}(\bar{F}))\rightarrow
R^{i}p_{Y^{\ast}}(\bar{h}^{\ast}(\bar{F}))$ pour $i=0,1$ est bijectif,
\resp pour tout faisceau abélien de $\LL$-torsion $\bar{F}$ sur
$\bar{X}$, l'homomorphisme de changement de base est bijectif pour
$i\leq n$).

\item[$3^{\circ}$)] L'image inverse de $\xi$ dans
$H^{i}(\bar{Y},p_Y\ast(h^{\ast}(F)))=H^{i}(\bar{Y},
\bar{h}^{\ast}(p^{\ast}(F)))$ est contenue dans l'image de
$H^{i}(\bar{X},p^{\ast}(F))$.
\end{itemize}

\item Soit\pageoriginale $n\geq -1$, et supposons que pour tout
faisceau d'ensembles $F$ (\resp tout faisceau de
ind-$\LL$-groupes, \resp tout faisceau abélien de
$\LL$-torsion) $F$ sur $X$, l'homomorphisme
$H^{i}(X,F)\rightarrow H^{i}(Y,h^{\ast}(F))$ soit bijectif pour $i\leq
n$, injectif pour $i=n+1$, (dans le premier cas, on suppose $n=1$,
dans le deuxième, $n=-1$ ou $0$). Soit $\xi \in
H^{n+1}(Y,h^{\ast}(F))$. Pour que $\xi$ soit dans l'image de
$H^{n+1}(X,F)$, il faut et il suffit qu'il satisfasse à la condition
énoncée dans {\rm(i)} ci-dessus, avec $i=n+1$).
\end{enumerate}
\end{enonce*}

\noindent
\emph{Démonstration de} \hyperref[XII.cor6.10]{6.10}. Utilisant \hyperref[XII.prop6.5]{6.5} 
et une récurrence
immédiate sur $n$, on constate que \hyperref[XII.cor6.10]{6.10} (i) est conséquence de 
\hyperref[XII.cor6.10]{6.10} (ii), que nous allons maintenant démontrer.

Supposons que $\xi$ ne soit pas dans l'image de $H^n(X,F)$. Soit
$\Phi$ l'ensemble des parties fermées $X'$ de $X$ telles que l'image
$\xi$ dans $H^n(Y',h^{\ast}(F)| Y')$ ne soit pas contenue dans
celle de $H^n(X',F| X')$. Par hypothèse, $X\in \Phi$, donc $\Phi$
n'est pas vide. Ordonnons $\Phi$ par la relation $\supset$, et
montrons que $\Phi$ est inductif. Pour ceci, il suffit de prouver que
si $(X'_{\lambda})$ est une famille totalement ordonnée d'éléments de
$\Phi$, et si $X'$ est leur intersection, alors $X'\in \Phi$. Or,
munissant les $X'_{\lambda}$ et $X'$ de la structure induite 
réduite, on voit que les morphismes $X'_{\lambda}\rightarrow
X$ sont affines (puisque ce sont des immersions fermées), donc le
système projectif des $X'_{\lambda}$ satisfait aux conditions
envisagées dans VII 5. De plus, on constate aussitôt que $X'$ est la
limite projective des $X'_{\lambda}$. De même, $Y'$ est la limite
projective du système projectif des $Y'_{\lambda}$, qui sont affines
sur $Y$. Utilisant VII 5.8, on trouve donc
$$
\varinjlim H^n(X'_{\lambda},F|X'_{\lambda})\xrightarrow{\sim}
H^n(X',F|X')\quav \varinjlim H^n(Y'_{\lambda},
h^{\ast}(F)|Y'_{\lambda})\xrightarrow{\sim}
H^n(Y',h^{\ast}(F)|Y'){\quoi}.
$$
De ceci, on conclut aussitôt que si on avait $(\xi|Y')\in \Iim
H^n(X',F|X')$, alors il\pageoriginale existerait un indice $\lambda$
tel que $(\xi|Y'_{\lambda})\in \Iim H^n(X'_{\lambda},F|X'_{\lambda})$, ce
qui est absurde; donc on a bien $X'\in \Phi$. Donc $\Phi$ est
inductif, et contient par suite un élément minimal, soit $X'$. Comme
évidemment $∅\not\in \Phi$, $X'$ est non-vide. Appliquons à $\xi$
et à $X'$ l'hypothèse de \hyperref[XII.cor6.10]{6.10}, d'où un morphisme $p:\bar{X}\rightarrow
X$ satisfaisant aux conditions $1^{\circ}$) à $3^{\circ}$) énoncées
dans \hyperref[XII.cor6.10]{6.10}. Notons qu'on peut supposer même $p$ \emph{surjectif}. En
effet, soit $U$ un ouvert non vide de $X'$ contenu dans $p(\bar{X})$,
et soit $X'_1$ son complémentaire dans $X'$, de sorte que par
construction on a $X'_1\not\in \Phi$. Utilisant cette relation, on
voit tout de suite que, posant $\bar{X}_1=\bar{X}\amalg X'_1$ et
désignant par $\bar{p}_1:\bar{X}_1\rightarrow X$ le morphisme défini
par $\bar{p}$ et l'inclusion de $X'_1$ dans $X$, le morphisme
$\bar{p}_1$ satisfait encore aux mêmes hypothèses que $\bar{p}$; de
plus, il est surjectif. Nous supposons donc $p$ surjectif, et
appliquons maintenant \hyperref[XII.cor6.9]{6.9}, en y remplaçant le morphisme
$h:Y\rightarrow X$ par le morphisme $h':Y'\rightarrow X'$ ($X'$ étant
muni, disons, de la structure induite réduite), et $\xi$ par
$\xi|Y'$. On conclut alors de \hyperref[XII.cor6.9]{6.9} que l'on a $(\xi|Y')\in \Iim
H^n(X',F|X')$, ce qui contredit la relation $X'\in \Phi$ et achève la
démonstration de \hyperref[XII.cor6.10]{6.10}.

Nous utiliserons \hyperref[XII.cor6.10]{6.10} pour ramener \hyperref[XII.cor5.5]{5.5} 
au cas où $f$ est un morphisme
projectif; mais pour ceci, dans le cas non noethérien, nous aurons
besoin d'une variante non noethérienne du lemme de Chow, qui sera
donnée au numéro suivant.

\begin{enonce*}{Proposition 6.11}\etiquette[6.11]{XII.prop6.11}
(Lemme de transitivité). Les notations sont celles de \hyperref[XII.prop6.8]{6.8}, mais on
ne fait aucune hypothèse de surjectivité relativement à $p$. On
suppose que $h$ satisfait à la condition {\rm (i) a)} (\resp {\rm (ii)
a)}, \resp {\rm (iii) a)} de \hyperref[XII.prop6.5]{6.5}. On suppose de plus, dans le cas
{\rm(i)}, que pour tout faisceau d'ensembles $\bar{F}$ sur $\bar{X}$,
l'homomorphisme de changement de base de \hyperref[XII.prop6.8]{6.8} est un monomorphisme
(\resp un isomorphisme); dans le cas {\rm(ii)}, que $\bar{X}$ est
noethérien et que pour tout faisceau $\bar{F}$ de
ind-$\LL$-groupes, l'homomorphisme de changement de base
$h^{\ast}(R^1p_{\ast}(\bar{F}))\rightarrow
R^1p_{Y^{\ast}}(\bar{h}^{\ast}(\bar{F}))$\pageoriginale est bijectif
pour $i=0,1$; enfin, dans le cas {\rm(iii)}, que l'homomorphisme de
changement de base est bijectif pour $i\leq n$, injectif pour
$i=n+1$. Alors le morphisme $\bar{h}:\bar{Y}\rightarrow \bar{X}$
satisfait également à la condition {\rm(i)} (\resp {\rm(ii)}, \resp
{\rm(iii)}) de \hyperref[XII.prop6.5]{6.5}.
\end{enonce*}

\noindent
\emph{Démonstration de} \hyperref[XII.prop6.11]{6.11}. Dans le cas (i), on considère le
diagramme commutatif de la démonstration de \hyperref[XII.prop6.8]{6.8} (i). Ici les
hypothèses impliquent que la première flèche horizontale et la
première flèche verticale de droite sont injectives (\resp
bijectives), donc il en est de même de la deuxième flèche horizontale,
ce qu'on voulait établir. Dans le cas (ii), on se donne un faisceau
$\bar{F}$ de ind-$\LL$-groupes sur $\bar{X}$, et il faut
montrer que tout élément de $H^1(\bar{Y},\bar{h}^{\ast}(\bar{F}))$
provient d'un élément de $H^1(\bar{X},\bar{F})$. En vertu de \hyperref[XII.cor6.6]{6.6} (ii)
appliqué à $\bar{h}$, il suffit de trouver un monomorphisme
$\bar{F}\rightarrow \bar{G}$ de faisceaux en groupes sur $\bar{X}$ qui
\emph{efface} $H^1(\bar{Y},\bar{h}^{\ast}(\bar{F}))$. Or considérons
la deuxième ligne du diagramme utilisé dans la démonstration de \hyperref[XII.prop6.8]{6.8}
(ii). Il suffit successivement de trouver un monomorphisme
$\bar{F}\rightarrow \bar{G}$, avec $\bar{G}$ un faisceau en groupes
in-$\LL$-finis, qui efface le dernier terme $H^{\circ}(Y,R^1
p_{Y^{\ast}}(\bar{h}^{\ast}(\bar{F}))$, puis un monomorphisme de
faisceaux en groupes $\bar{G}\rightarrow \bar{H}$ qui efface
$H^1(Y,p_{Y^{\ast}}(\bar{h}^{\ast}(\bar{F}))$. Or, comme on a remarqué
dans la démonstration de \hyperref[XII.prop6.8]{6.8} (i), les deux termes qu'il s'agit
d'effacer sont isomorphes respectivement, par les flèches verticales
extrêmes du diagramme envisagé, à $H^{\circ}(X,R^1p_{\ast}(\bar{F}))$
et à $H^1(X,p_{\ast}(\bar{F}))$, compte tenu de l'hypothèse sur les
homomorphismes de changement de base faites ici sur $p$,
$h$. D'ailleurs, les isomorphismes envisagés sont évidemment
fonctoriels en $\bar{F}$, de sorte qu'il suffit d'effacer les deux
termes précédents. Or en vertu de \hyperref[XII.prop3.3]{3.3} on peut effacer
$R^1p_{\ast}(\bar{F})$ (et à fortiori
$H^{\circ}(X,R^1p_{\ast}(\bar{F}))$ par un monomorphisme de $\bar{F}$
dans un ind-$\LL$-groupe $\GG$. Il reste à effacer
$H^1(X,p_{\ast}(\bar{F}))$, ce qui est possible grâce au fait que par
le\pageoriginale diagramme envisagé, ce dernier terme s'envoie dans
$H^1(\bar{X},\bar{F})$ par un monomorphisme fonctoriel en $\bar{F}$,
et que $H^1(\bar{X},\bar{F})$ est encore effaçable grâce à \hyperref[XII.prop3.3]{3.3}.

Il reste à traiter le cas (iii), et pour celui-ci encore nous
reprenons à rebours la démonstration de \hyperref[XII.prop6.8]{6.8} \footnote[1]{On pourrait
aussi, au lieu de \hyperref[XII.cor6.3]{6.3}, utiliser la suite spectrale de Leray pour $p$,
$p_Y$.}. Il faut prouver que l'homomorphisme
$$
H^{i}(\bar{X},\bar{F})\longrightarrow
H^{i}(\bar{Y},\bar{h}^{\ast}(\bar{F}))
$$
est un isomorphisme pour $i\leq n$, et un monomorphisme pour
$i=n+1$. Or les homomorphismes précédents sont induits par un
homomorphisme de complexes
$$
\RR\Gamma_{\bar{X}}(\bar{F})\longrightarrow
\RR\Gamma_{\bar{Y}}(\bar{h}^{\ast}(\bar{F})){\quoi},
$$
déduit d'une résolution injective $C(\bar{F})$ et d'un homomorphisme
de $\bar{h}^{\ast}(C(\bar{F}))$ (qui est une résolution de
$h^{\ast}(\bar{F})$) dans une résolution injective de
$\bar{h}^{\ast}(\bar{F})$. L'homomorphisme ci-dessus s'identifie
également à l'homomorphisme composé
$$
\RR\Gamma_X(Rp_{\ast}(\bar{F}))\longrightarrow
\RR\Gamma_Y(h^{\ast}(\bar{F}))\longrightarrow
\RR\Gamma_Y(Rp_{Y^{\ast}}(\bar{h}^{\ast}(\bar{F})){\quoi},
$$
où le premier homomorphisme est l'homomorphisme $h^{\ast}$ relatif à
$\RR p_{\ast}(\bar{F})$, et le deuxième est déduit des
homomorphismes de changement de base
$h^{\ast}(\RR p_{\ast}(\bar{F}))\rightarrow
\RR p_{Y^{\ast}}(\bar{h}^{\ast}(\bar{F}))$ en appliquant
$\RR\Gamma_Y$. Le premier homomorphisme induit donc un
isomorphisme sur les groupes de cohomologie des complexes envisagés,
en vertu de l'hypothèse faite sur $h$ et de \hyperref[XII.cor6.3]{6.3} (i). Il en est de même
pour le deuxième, en vertu de l'hypothèse que nous avons faite sur les
homomorphismes de changement de base, en utilisant l'argument donné
dans la démonstration de \hyperref[XII.prop6.8]{6.8} (iii). Par suite, la même conclusion
s'applique au composé des deux, ce qui achève de prouver \hyperref[XII.prop6.11]{6.11}.

\begin{enonce*}[remark]{Remarque 6.12}\etiquette[6.12]{XII.rem6.12}
L'argument\pageoriginale donné pour \hyperref[XII.prop6.11]{6.11} dans les cas (i) et (iii)
donne encore un résultat plus précis: pour pouvoir conclure, pour un
faisceau $\bar{F}$ \emph{donné} sur $\bar{X}$, que
$H^{i}(\bar{X},\bar{F})\rightarrow
H^{i}(\bar{Y},\bar{h}^{\ast}(\bar{F}))$ est bijectif pour $i\leq n$,
injectif pour $i=n+1$, il suffit de supposer
\begin{itemize}
\item[$1^{\circ}$)] l'homomorphisme de changement de base
$h^{\ast}(R^{i}p_{\ast}(\bar{F}))\rightarrow
R^{i}p_{Y^{\ast}}(\bar{h}^{\ast}(\bar{F}))$ est un isomorphisme pour
$i\leq n$, un monomorphisme pour $i=n+1$, et

\item[$2^{\circ}$)] pour tout couple d'entiers $i$, $j\geq 0$,
l'homomorphisme $H^{i}(X,\RR^jp_{\ast}(\bar{F}))\rightarrow
H^{i}(Y,h^{\ast}(\RR^jp_{\ast}(\bar{F}))$ induit par $h$ est un
isomorphisme si $i+j\leq n$, un monomorphisme si $i+j=n+1$.
\end{itemize}

Il est plausible qu'on a un résultat analogue dans le cas (ii), mais
son énoncé devrait nécessairement faire intervenir de la 2-cohomologie
non commutative. C'est pour éviter le recours à cette théorie que nous
avons utilisé un argument différent dans \hyperref[XII.cor6.10]{6.10} (ii), utilisant les
propriétés d'effaçabilité du $H^1$ non commutatif, argument qui ne
donne pas en revanche de résultat, comme ci-dessus, pour un $\bar{F}$ fixé.
\end{enonce*}

\begin{enonce*}[remark]{Remarques 6.13}\etiquette[6.13]{XII.rems6.13}
Le cas le plus important où les conditions de \hyperref[XII.prop6.5]{6.5} (i) (ii) (iii) sont
remplies est celui envisagé dans \hyperref[XII.cor5.5]{5.5}, qui sera prouvé dans l'exposé
suivant en utilisant les réductions faites dans le présent exposé. Un
autre cas intéressant (où $h$ n'est plus une immersion fermée comme
dans \hyperref[XII.cor5.5]{5.5} sera étudié dans $XV$. Signalons également le cas, assez voisin
de 5.5, où $X$ est le spectre d'un anneau noethérien, séparé et
complet pour la topologie définie par un idéal $J$, et où $Y=\Spec
(A/J)$ est le sous-schéma fermé de $X$ défini par $J$. On vérifie
alors directement les conditions de \hyperref[XII.prop6.5]{6.5} (ii) (avec
$\LL=\PP$) sous la forme\pageoriginale ${\rm b''}$) (\EGA
IV 18.3.2); il est très plausible que les conditions de \hyperref[XII.prop6.5]{6.5} (iii) sont
également vérifiées pour tout $n$, \ie que pour tout faisceau de
torsion $F$ sur $X$, les homomorphismes $H^n(X,F)\rightarrow
H^n(Y,F|Y)$ soient bijectifs. On peut faire des conjectures
considérablement plus générales, liées à une généralisation à des
schémas non locaux de la construction de hensélisation. Tout d'abord,
notons qu'on ne connaît pas d'exemple, $h$ étant une immersion fermée,
où les conditions de \hyperref[XII.prop6.5]{6.5} (ii) (iii) ne soient satisfaites (avec
$\LL=\PP$) lorsque la condition respée de \hyperref[XII.prop6.5]{6.5} (i) l'est,
c'est-à-dire lorsque le couple $(X,Y)$ est un \emph{couple hensélien}
dans la terminologie de (\EGA IV 18.5.5); on peut donc conjecturer
qu'un tel couple satisfait aux conditions de \hyperref[XII.prop6.5]{6.5} (ii) (iii) (avec
$\LL=\PP$, tout $n$). D'autre part, on peut se demander
sous quelles conditions la propriété suivante pour un couple $(X,Y)$
d'un schéma $X$ quasi-complet et quasi-séparé et d'un sous-schéma
fermé $Y$ implique que le couple est hensélien:

$(\ast)$ Pour tout schéma $X'$ étale sur $X$, posant
$Y'=\fprod{X'}{Y}{X}$, l'application canonique
$\Gamma(Y'/Y)\rightarrow \Gamma(X'/X)$ est bijective.
(\Cf \EGA IV 18.5.4 b)). Cette condition est évidemment plus faible à
priori que la condition hensélienne \hyperref[XII.prop6.5]{6.5} (i); elle est même strictement
plus faible, car elle est vérifiée par exemple si $X$ est un schéma
projectif normal irréductible de dimension $\geq 2$ sur un corps $k$,
et $Y$ une section hyperplane de $X$ (comme il résulte facilement de
\SGA 2 XII 2.1, où on fait $S=\Spec(k)$), mais évidemment le couple
$(X,Y)$ n'est pas hensélien. Il est possible par contre que si $X$ est
affine, la condition $(\ast)$ implique que $(X,Y)$ est un couple
hensélien; s'il en était ainsi, pour tout schéma affine $X$ et tout
partie fermée $Y$ de $X$, un\pageoriginale procédé de localisation
étale le long dé $Y$, analogue à celui de la hensélisation des anneaux
locaux, devrait permettre de lui associer un couple \emph{hensélien}
$(X',Y)$, où $X'$ est « proétale »  sur $X$. Une autre hypothèse
sur $(X,Y)$ (suggérée par des résultats de Hironaka et Rossi sur la
contractibilité de certaines sous-variétés de variétés algébriques),
qui permettrait peut-être de déduire la propriété hensélienne de la
propriété plus faible $(\ast)$, est la suivante: l'immersion
$Y\rightarrow X$ est régulière de codimension 1, et le faisceau
conormal $\NN_{Y/X}=\JJ/\JJ^2$ (où $\JJ$ est l'idéal qui définit $Y$
dans $X$) est un faisceau inversible \emph{ample} sur $Y$. (NB. Dans
le cas du contre-exemple envisagé plus haut, ce faisceau conormal
était un contraire anti-ample, \ie son inverse étant ample).
\end{enonce*}

\begin{NoHyper}
\footnotetext{Le lecteur qui ne s'intéresse aux énoncés du par. 5
que dans le cas noethérien peut omettre la lecture du présent numéro.}
\end{NoHyper}
\section{Une variante du Lemme de Chow ${}^{(*)}$.}\etiquette[7]{XII.sec7}

Nous aurons besoin de la variante suivante de (\EGA II 5.6.1)
débarrassée de toutes hypothèses noethérienne ou d'irréductibilité.

\begin{enonce*}{Lemme 7.1}\etiquette[7.1]{XII.lem7.1}
Soit $S$ un schéma quasi-compact et quasi-séparé, et $f:X\rightarrow
S$ un morphisme séparé de type fini, avec $X$ non vide. Il existe un
diagramme commutatif
$$
\xymatrix@=1.25cm{
X \ar[rd]_{f}&    &\bar{X}\ar[ll]_{\pi}\ar[ld]^{\bar{f}}\\
  &  S & 
}
$$
avec\pageoriginale $\pi$ projectif et $\bar{f}$ quasi-projectif, et un
ouvert non-vide $U$ de $X$, tels que $\pi$ induise un isomorphisme de
$\bar{U}=\fprod{U}{\bar{X}}{X}=\pi^{-1}(U)$ sur $U$.
\end{enonce*}

\noindent
\emph{Démonstration}. Notons d'abord: Soit $i:X'\rightarrow X$ un
sous-schéma fermé de $X/S$ et $V\subset X'$ un sous-schéma ouvert dans
$X$ et dense dans $X'$. Si le lemme est vrai pour $X'$, il l'est pour
$X$. En effet, soit
$$
\xymatrix@=1.25cm{
X' \ar[rd]_{f'}&    &\bar{X}'\ar[ll]_{\pi'}\ar[ld]^{\bar{f}'}\\
  &  S & 
}
$$
un diagramme comme dans l'énoncé, et $U'\neq \phi$ un ouvert de $X'$
tel que $\pi'$ induise un isomorphisme
$\bar{U}'={\pi'}^{-1}(U')\rightarrow U'$. Puisque $V$ est dense dans
$X'$, $U'\cap V\neq \phi$ et on peut donc remplacer $U'$ par cet
ouvert. Alors $U'$ est un ouvert de $X$, et on pose $\pi=i \pi'$,
$\bar{f}=\bar{f}'$, $U=U'$.

Or il existe un recouvrement fini de $X$ par des ouverts affines non
vides $U_k$, $k=1,\ldots,n$; de l'hypothèse sur $S$ résulte que chaque
$U_k\rightarrow S$ est quasi-affine, et a fortiori
quasi-compact. Démontrons le théorème pour tout $X$, par récurrence
sur $n$: Si l'intersection $U$ des $U_k$ est dense dans $X$ on peut
copier la démonstration de (\EGA II 5.6, B, C, D), tenant compte des
modifications de la notation. Si l'intersection est non-vide, on se
ramène au cas où elle est dense en remplaçant $X$ par l'adhérence
$X'$ de $U$ (\cf plus haut), et $U_k$ par $U'_k=U_k\cap X'$. Supposons
enfin que $U$ soit vide et choisissons $r$, $1\leq r\leq n$, tel que
$V=\cap^n_{k=1}U_k\neq \phi$ mais $V\cap U_{r+1}=\phi$ - c'est
évidemment possible puisque $U_1\neq \phi$. Nous pouvons maintenant
remplacer $X$ par l'adhérence $X'$ de $V$ et les $U$ par les $U_k\cap
V$ non-vides. Comme alors $U_{r+1}\cap V$ est vide, on obtient le
résultat par l'hypothèse de récurrence.

\section{Réductions définitives}\etiquette[8]{XII.sec8}\pageoriginale

\begin{enonce*}{Lemme 8.1}\etiquette[8.1]{XII.lem8.1}
Pour prouver l'un des énoncés {\rm \hyperref[XII.thm5.1]{5.1} (i) (ii) (iii)}, il suffit de
le faire lorsque l'on suppose $f$ \emph{projectif} et $S$ noethérien.
\end{enonce*}

Supposons d'abord \hyperref[XII.thm5.1]{5.1} démontré pour $f$ projectif, montrons comment on
en conclut le même énoncé pour tout $f$. Utilisant \hyperref[XII.lem6.1]{6.1}, on est réduit
au cas de \hyperref[XII.cor5.5]{5.5}, avec $S$ strictement local. Appliquons maintenant \hyperref[XII.cor6.10]{6.10}
en faisant $Y=X_{\circ}$. Étant donné $(\xi, X')$ comme énoncé dans
\hyperref[XII.cor6.10]{6.10}, nous appliquons le lemme de Chow \hyperref[XII.lem7.1]{7.1} pour trouver un
$S$-morphisme $p:\bar{X}\rightarrow X'$, avec $\bar{X}$ projectif sur
$S$, tel que $p(\bar{X})$ contienne un ouvert non vide de $X'$
(condition $1^{\circ}$) de \hyperref[XII.cor6.10]{6.10}). Comme $p$ est projectif, par
hypothèse les homomorphismes de changement de base pour $(p,h)$ sont
des isomorphismes (condition $2^{\circ}$) de \hyperref[XII.cor6.10]{6.10}). Enfin, comme $X$
est projectif sur $S$, par hypothèse le morphisme d'inclusion
$\bar{Y}=\bar{X}_{\circ}\rightarrow \bar{X}$ induit un isomorphisme
pour les groupes de cohomologie $H^{i}(\bar{X},\bar{F})\rightarrow
H^{i}(\bar{X}_{\circ},\bar{h}^{\ast}(\bar{F}))$, (condition
$3^{\circ}$) de \hyperref[XII.cor6.10]{6.10}). Donc les conditions de \hyperref[XII.cor6.10]{6.10} sont bien
satisfaites, d'où la conclusion.

Montrons maintenant que pour vérifier l'un des trois énoncés \hyperref[XII.thm5.1]{5.1} pour
toute situation $(f,F,g)$, avec $f$ projectif, on peut supposer de
plus que $S$ est noethérien. Pour ceci, on note qu'on peut supposer
grâce à \hyperref[XII.lem6.1]{6.1} que $S$ est strictement local et que $S'\rightarrow S$ est
l'inclusion du point fermé, et on est ramené à vérifier l'énoncé \hyperref[XII.cor5.5]{5.5}
avec $S$ strictement local. De plus, comme $X$ se plonge dans le
schéma projectif-type $\PP^r_S$, on peut (grâce à la compatibilité \hyperref[XII.prop4.4]{4.4}
(ii)) remplacer $X$ par $\PP^r_S$, $F$ par son image dans $\PP^r_S$,
ce qui nous ramène au cas où $X=\PP^r_S$. Notons que $A$ est limite
inductive filtrante de sous-anneaux\pageoriginale locaux henséliens
noethériens $A_i$ tels que l'homomorphisme d'inclusion $A_i\rightarrow
A$ soit local: prendre les hensélisés stricts des sous-$\ZZ$-algèbres
de type fini de $A$ en les idéaux premiers induits par l'idéal maximal
de $A$. Sous ces conditions, le corps résiduel $k$ de $A$ est
également la limite inductive des corps $k$ des $A_i$ (\EGA IV
5.13.1). Soient $S=\Spec (A_i)$, $X_i=\PP^r_{S_i}$, et supposons
donné un système inductif $(F_i)$ de faisceaux d'ensembles
(\resp\ldots) sur les $X_i$ comme dans VII 5.7, et soit $F$ le
faisceau qu'il définit sur $X$. Soit $F_{i\circ}$ (\resp $F_{\circ}$)
le faisceau induit par $F_i$ sur la fibre $X_{i\circ}$ du point fermé
$s_i$ de $S_i$ (\resp sur la fibre $X_{\circ}$ du point fermé $s$ de
$S$). Alors on a un diagramme commutatif
$$
\xymatrix{
\displaystyle\mathop{\varinjlim}_{i} H^{n}(X_i,F_i) \ar[r] \ar[d]&
H^{n}(X,F)\ar[d]\\ 
\displaystyle\mathop{\varinjlim}_{i} H^{n}(X_{i0},F_{i0}) \ar[r] &
H^{n}(X_0,F_0), 
}
$$
dont les deux flèches horizontales sont des isomorphismes, en vertu de
VII 5.7. Par suite, pour prouver que la deuxième flèche verticale est
un isomorphe, il suffit de le prouver pour la première, donc il suffit
de prouver que les homomorphismes $H^n(X_i,F_i)\rightarrow
H^n(X_{i\circ},F_{i\circ})$ sont des isomorphismes. D'autre part, tout
faisceau d'ensembles (\resp\ldots) $F$ sur $X$ est isomorphe à la
limite d'un système inductif comme ci-dessus, en prenant par exemple
$F_i=u_{i^{\ast}}(F)$, $u_i:X\rightarrow X_i$ étant le morphisme
canonique, en vertu de (\hyperref[XII.cor5.2]{5.2}) et (\hyperrefext[IX][IX.prop1.6]{1.6}
(iii)). Cela achève la 
démonstration de \hyperref[XII.lem8.1]{8.1}.

\begin{enonce*}{Corollaire 8.2}\etiquette[8.2]{XII.cor8.2}
Le théorème {\rm 5.1 (i)} est vrai.
\end{enonce*}

En effet, il suffit de conjuguer \hyperref[XII.lem8.1]{8.1}, le critère \hyperref[XII.prop6.5]{6.5} (i) ${\rm b''}$), et \hyperref[XII.lem5.8]{5.8}.

\begin{enonce*}{Lemme 8.3}\etiquette[8.3]{XII.lem8.3}
Pour\pageoriginale prouver l'un des énoncés {\rm \hyperref[XII.thm5.1]{5.1} (i) (ii) (iii)},
il suffit de le faire lorsqu'on suppose $f$ projectif de dimension
relative $\leq 1$, et $S$ noethérien.
\end{enonce*}

En effet, nous savons déjà par \hyperref[XII.lem8.1]{8.1} qu'on peut se borner au cas $f$
projectif, $S$ noethérien, et un argument déjà signalé nous permet de
supposer de plus qu'on a $X=\PP^r_S$. Nous procédons alors par
récurrence sur $r$, en notant que le théorème est vrai par hypothèse
pour $r=1$. On peut donc supposer $r\geq 2$, et le théorème démontré
pour $\PP^{r-1}_S\rightarrow S$ (toute base $S$). Considérons
$Z=\PP^{r-2}_S$ comme plongé dans $X=\PP^r_S$ de la façon
habituelle, et soit $p:\bar{X}\rightarrow X$ obtenu en faisant éclater
$Z$ dans $X$. Un calcul immédiat (\cf \EGA V) montre que l'on a un
morphisme naturel $p_1:\bar{X}\rightarrow X_1=\PP^1_S$, qui fait de
$\bar{X}$ un $X_1$-schéma localement isomorphe à $\PP^{r-1}_{X_1}$
(c'est le fibré projectif sur $X_1$ associé à un faisceau localement
libre de rang $r$ convenable sur $X_1$). D'autre part, les fibres de
$p$ sont de dimension $\leq 1$, ainsi que celles de
$f_1:X_1\rightarrow X$. Nous pouvons donc appliquer l'hypothèse, et
l'hypothèse de récurrence, aux trois morphismes $p$, $p_1$, $f_1$ dans
le diagramme commutatif
$$
\xymatrix{
   & \bar{X} \ar[ld]_{p} \ar[rd]^{p_1} &          \\
 X \ar[rd]_{f}&           & X_1\ar[ld]^{f_1}\\
   &    S      &       
}$$
Comme $p$ est surjectif, le lemme \hyperref[XII.prop6.8]{6.8} nous ramène à prouver \hyperref[XII.thm5.1]{5.1} pour
le composé $\bar{f}=fp=f_1p_1$. Or nous pouvons déjà appliquer \hyperref[XII.thm5.1]{5.1} aux
deux morphismes $f_1$ et $p_1$. Compte tenu de \hyperref[XII.prop6.11]{6.11}, on en conclut que
\hyperref[XII.thm5.1]{5.1} est vrai également pour leur composé $f_1p_1$. Cela achève la
démonstration de \hyperref[XII.lem8.3]{8.3}.

\begin{enonce*}{Lemme 8.4}\etiquette[8.4]{XII.lem8.4}
Pour\pageoriginale démontrer {\rm \hyperref[XII.thm5.1]{5.1} (ii)}, il suffit de
démontrer le 
« théorème de spécialisation pour le groupe fondamental » {\rm
\hyperref[XII.thmbis5.9]{5.9 bis}}, dans le cas où $f$ est projectif de
dimension relative $\leq 
  1$, et où $S$ est strictement local noethérien.
\end{enonce*}

Cela résulte en effet aussitôt de \hyperref[XII.lem8.3]{8.3}, \hyperref[XII.lem6.1]{6.1}
et \hyperref[XII.prop6.5]{6.5} (ii) (critère b)). 

On démontrera directement le théorème de spécialisation, dans le cas
\hyperref[XII.lem8.4]{8.4}, dans l'exposé suivant (\hyperrefext[XIII][XIII.sec2]{2}).

\begin{enonce*}{Lemme 8.5}\etiquette[8.5]{XII.lem8.5}
Pour démontrer {\rm \hyperref[XII.thm5.1]{5.1} (iii)}, il suffit de démontrer
{\rm \hyperref[XII.thm5.1]{5.1} (ii)} 
et, de plus, la proposition suivante (qui sera démontrée dans l'exposé
suivant {\rm (\hyperrefext[XIII][XIII.sec3]{3})} ou {\rm(\EGA IV 21.9.12)}):
\end{enonce*}

\begin{enonce*}{Proposition 8.6}\etiquette[8.6]{XII.prop8.6}
Soit $S$ hensélien et noethérien, $f:X\rightarrow S$ projectif de
dimension relative $\leq 1$, et $X_{\circ}$ la fibre fermée de
$X/S$. Pour tout schéma $Z$, soit $\Pic Z=H^1(Z,(G_m)_Z)$ le groupe de
Picard de $Z$. Le morphisme de restriction
$$
\Pic X\longrightarrow \Pic X_{\circ}
$$
est surjectif.
\end{enonce*}

\noindent
\emph{Démonstration de} \hyperref[XII.lem8.5]{8.5}. D'après \hyperref[XII.prop6.5]{6.5} 
(iii) et \hyperref[XII.cor8.2]{8.2}, il suffit de démontrer que
$$
\varphi^q:H^q(X,\ZZ/\ell^\nu \ZZ)\longrightarrow
H^q(X_{\circ},\ZZ/\ell^{\nu}\ZZ)
$$
est surjectif pour tout $q\geq 1$, et tout nombre premier 1. Or on
connaît la surjectivité si $q=1$, d'après \hyperref[XII.thm5.1]{5.1} (ii). De
plus, puisque 
$X_{\circ}$ est un schéma projectif de dimension $\leq 1$ sur un corps
$k$ séparablement clos, on a
$$
H^q(X_{\circ},\ZZ/\ell^{\nu}\ZZ)=0
$$
si $q>2$ dans le cas $\ell\neq \car. k$ (\hyperrefext[X][X.cor4.3]{4.3}), et si $q>1$ dans le
cas $\ell=\car. k$ (\hyperrefext[X][X.cor5.2]{5.2}). La surjectivité est donc triviale pour ces
valeurs de $q$! 

Il\pageoriginale reste à traiter le cas $q=2$ et $\ell\neq
\car.k$. Notons $n=\ell^{\nu}$. Comme le faisceau des racines
$n$-ièmes de l'unité $(𝛍_n)_S$ est localement
isomorphe à $(\ZZ/n)_S$, et que $S$ est strictement local, il
existe un isomorphisme (non canonique)
$(𝛍_n)_S\simeq (\ZZ/n\ZZ)_S$ d'où
$(𝛍_n)_X\simeq (\ZZ/n\ZZ)_X$.

Par théorie de Kummer \hyperrefext[IX][IX.Thr3.2]{3.2} appliquée à $X$ et $X_{\circ}$, on a un
morphisme de suites exactes
$$
\xymatrix@R=1.5cm{
\Pic X \ar[r] \ar[d]^{\epsilon}& H^{2}(X, {\phantom{a}}_n) \ar[d]^{\varphi^2}\ar[r]& H^{2}(X,G_m)\ar[d]\\
\Pic X_0 \ar[r]& H^2(X_0, {\phantom{a}}_n) \ar[r]& H^2(X_0,G_m).
}
$$
Or $H^2(X_{\circ},G_m)=0$ d'après \hyperrefext[IX][IX.thm4.6]{4.6}. Donc la surjectivité de
$\epsilon$ implique celle de $\varphi^2$, d'où le lemme.

\begin{thebibliography}{99}
\bibitem{SGA4-12:1} \textsc{Giraud} J., Cohomologie non-abélienne, thèse (à
paraître), et deux notes aux C.R. t.260, p.2392 et 2656 (1965).

\bibitem{SGA4-12:2} \textsc{Godement} R., Théorie des faisceaux, Paris (1958).

\bibitem{SGA4-12:3} \textsc{Verdier} J.L., Catégories dérivées, Institut des
Hautes Études Scientifiques (1964).
\end{thebibliography}

\ifx\danslelivre\undefined
\end{document}
\fi
